
\chapter[Multiple scattering formulations]
{Multiple scattering formulations for a stratified half-space 
with superimposed heterogeneity} \label{stratG}








In this chapter, we derive
representations of the {\em elastodynamic Lippmann-Schwinger\/} 
equation (Lippmann, \& Schwinger, 1950), appropriate to a stratified halfspace 
with superimposed lateral heterogeneity. When we attempt to solve this 
equation in terms of the solutions to the component subregions (usually, 
either thin slabs, or small squares) of the layer stack, we are effectively developing a {\em multiple scattering 
theory\/}.
Multiple scattering theory goes back at least as far as {\em Rayleigh\/}~(1892)
and has had numerous contributors since that time. Mathematically 
rigorous formulations were given by {\em Foldy\/}~(1945)
and {\em Lax\/}~(1951), and their development proceeds from a pair of 
axioms, referred to here as the {\em Multiple Scattering Axioms\/}.
Although the axioms are physically plausible, they are  not self-evident.

 Multiple scattering theory for a stratified half-space may be 
approached via two complementary routes:
\begin{itemize}
\item a generalization of the {\em numerical analysis\/} approach of {\em Chin, Hedstrom \& 
Thigpen\/}~(1984). While the standard method is based on a set of 
displacement-traction balances for the interfaces between each plane 
layer, our method is derived from a set of {\em wavefield balances\/}. 
The wavefield balance for a subregion is a statement that the scattered 
field is related to the incident field through the agency of the 
subregion scattering operator
\item the {\em formal scattering theory\/} approach, [e.g. {\em 
Kennett\/}~(1981), {\em Newton\/}~(1982)]. This approach is more abstract 
(and consequently, more general) 
than the numerical analysis route, but this is not necessarily a 
hindrance, as it allows us to to translate some of the theory developed in 
quantum mechanics to the present elastic wave theory. This approach is 
particularly useful for the investigation of surface waves. 
\end{itemize}
We will have cause to investigate both routes to multiple scattering 
theory. In \S\ref{RTfactor} we investigate the numerical analysis method, 
while the formal scattering method is presented in \S\ref{sectionstratG}.

In chapter \ref{Grep}, we introduced a {\em secondary source lattice\/} 
comprised of $L-$layers, with $N_x$ source points per layer. The entire 
lattice was 
embedded in free space. 
In this chapter, {\em each row of the lattice, is 
embedded in a separate plane layer. 
\/}
To map the wavefield between the lattice points,
we  generalize the planewave  propagator 
representations of chapter \ref{Grep} to a stratified halfspace.

Let us assume that we have a plane-stratified halfspace spanning the depth 
interval $0 \le z \le \infty$ (Fig.\ref{BasicStack}), with superimposed 
lateral heterogeneity supported on the compact depth interval 
$0<z_{0} \le 
z \le z_{L}<\infty$.
For each of the plane layers in the depth interval $z_{0} \le 
z \le z_{L}$, 
the material parameters ${\bmi  m}=(\rho,\alpha,\beta)$ are split into the 
sum of a homogeneous component
\[
{\bmi  m}_{i}\equiv  \left({ \rho }_{i},{ \alpha }_{i},{ \beta }_{i}
\right):i=0,1,\cdots ,L+1
\]
and a laterally heterogeneous component 
\[
{\mit \Delta}{\bmi  m}_{i}(x)\equiv \left({\mit \Delta}{ \rho }_{i}(x),
{\mit \Delta}{\alpha }_{i}(x),{\mit \Delta}{\beta}_{i}(x)\right):i=1,\cdots ,L
\]
({\em cf.\/}Fig.\ref{BasicStack}). The laterally heterogeneous component 
${\mit \Delta}{\bmi  m}_{i}$ in layer-$i$ is split 
into the sum of $N_x$ small squares, i.e.  
\[
{\mit \Delta}{\bmi  m}_{i}(x)\equiv \{ {\mit \Delta}{\bmi  m}_{i,j}: 1 
\le j \le N_x \}
\]


Note that the heterogeneity is confined to layers $i=1,\cdots ,L$, while 
the free surface, zone $0=z_{f}\le z \le z_{0}$ and the lower halfspace $z>z_{L}$
remain unperturbed in their plane stratified state. Also, we have 
assumed that the thickness ${\mit \Delta}z$ of each reference layer
$i=1,\cdots ,L$, is small on the scale of the incident wavelength, so 
that vertical variations of the material heterogeneity may be ignored. 
Consequently, 
the perturbations in the material parameters ${\mit \Delta}{\bmi  m}_{i}$ 
superimposed on each reference layer, 
are functions of the offset $x$ only, i.e. ${\mit \Delta}{\bmi  m}_{i}(x)$. This 
assumption is not essential for the validity of the theory, but simplifies the development.

\begin{figure}
\HideDisplacementBoxes
%\ShowDisplacementBoxes
\centerline{\BoxedEPSF{BasicStack.EPSF scaled 1000}}
\caption{The  labelling convention employed to model a stratified 
halfspace with superimposed heterogeneity}
\protect\label{BasicStack}
\end{figure}



We are just about ready to launch into the central part of the thesis, 
were we attempt to {\em build the reflection and transmission response of 
a stratified medium with superimposed heterogeneity, directly in terms of 
the reflection and transmission operators of the component subregions\/}. 
At this juncture, we pause to define our reflection/transmission operator 
notation [which follows {\em Bostrom \& Karlson\/}~(1984)]. 
This notation  is slightly amended from that adopted in {\em Kennett\/}~(1983). 
Throughout this work,  superscripts $(\cdot,\cdot)$ on scattering
operators indicate their spatial domains of dependence. The subscript arrows,
indicate the sense of the incident and scattered radiation {\em read right to left, in
typical matrix fashion}, i.e. the right arrow  indicates the sense of the incident field, and
the left arrow, the sense of the scattered field.
This notation is best explained by a specific example.
Recalling the mathematical notation 
\[
(z_1,z_2)\equiv z_{1} < z < z_{2}, \quad \mbox{and} \quad
[z_1,z_2]\equiv z_{1} \le z \le z_{2} 
\] 
and with reference to Fig.\ref{Sfig}, we write:
\begin{itemize}
\item ${\bmi S}_{\uparrow\downarrow }^{(1 , 2)}\left({\bmi 
S}_{\uparrow\downarrow }^{[1 , 2]}\right)$
for the reference medium downward incidence, upward reflection operator for the 
{\em open\/}({\em closed\/}) depth interval 
$z_1 <   z <  z_2\, (z_1 \le   z \le  z_2)$, which is written as ${\bmi R}_{D }^{1  2}$ in the 
notation of {\em Kennett\/}~(1983)
\item ${\bmi S}_{\uparrow\uparrow }^{(1 , 2)}\left({\bmi S}_{\uparrow\uparrow }^{[1 , 2]}\right)$
for the reference medium upward incidence, upward transmission operator for the 
{\em open\/}({\em closed\/}) depth interval 
$z_1 <   z <  z_2\, (z_1 \le   z \le  z_2)$, which is written as ${\bmi T}_{U }^{1  2}$ in the 
notation of {\em Kennett\/}~(1983)
\end{itemize}
Hence, ${\bmi S}_{\uparrow\downarrow }^{(1,2)}$ maps   $\downarrow\!  \!-going$
$P$,$SV$  and $SH$ plane wave components at depth $ {z}_{1} $ to the  $\uparrow\!  \!-going$ 
$P$,$SV$  and $SH$ plane wave components at the same level. Similarly, ${\bmi 
S}_{\uparrow\downarrow }^{[1 , 2]}$
is the reference medium downward reflection operator for the {\em closed\/} depth interval 
$z_1 \le   z \le  z_2$. The distinction between {\em open\/} and {\em 
closed\/} depth intervals is important when the material properties are discontinuous at 
the end points of the depth interval. While the open interval excludes 
the discontinuous change in the material properties, the closed interval 
includes them. Recalling the {\em adopted ordering\/} of the wavestates 
in $6 N_k \times 6 N_k$ array ${\bmi D}_{z}({\bf k}_{z},k_{y})$
in (\ref{Dsuperdef}),the ${\bmi S}_{\uparrow\downarrow }^{(1 , 2)}$ matrix 
has the component representation
\begin{equation} \label{Ssparsity}
{\bmi S}{}_{\uparrow\downarrow }^{(1,2)} \equiv 
\left({\begin{array}{ccc}
({\bmi S}{}_{\uparrow\downarrow }^{(1,2)})_{  PP}&
({\bmi S}{}_{\uparrow\downarrow }^{(1,2)})_{  PS}&
({\bmi S}{}_{\uparrow\downarrow }^{(1,2)})_{  PH}\\ 
({\bmi S}{}_{\uparrow\downarrow }^{(1,2)})_{  SP}&
({\bmi S}{}_{\uparrow\downarrow }^{(1,2)})_{  SS}&
({\bmi S}{}_{\uparrow\downarrow }^{(1,2)})_{  SH}\\ 
({\bmi S}{}_{\uparrow\downarrow }^{(1,2)})_{  HP}&
({\bmi S}{}_{\uparrow\downarrow }^{(1,2)})_{  HS}&
({\bmi S}{}_{\uparrow\downarrow }^{(1,2)})_{  HH} 
\end{array}}\right) 
\end{equation} 
where the $N_k \times N_k$ downward incident $c$-wave 
to upward reflected $c^{\prime}$-wave array 
${\left({\bmi S}{}_{\uparrow\downarrow 
}^{(1,2)}\right)}_{c^{\prime},c}$  is defined as
\begin{eqnarray}\label{RdspDef}
{\left({\bmi S}{}_{\uparrow\downarrow }^{(1,2)}\right)}_{c^{\prime},c}&=&
\mbox{diag}\left[\begin{array}{cccc}
\scriptstyle (-, c^{\prime} ,0) \leftarrow (+, c ,0),&
\scriptstyle (-, c^{\prime} ,1) \leftarrow (+, c ,1),&
\cdots, &
\scriptstyle (-, c^{\prime} ,M-1) \leftarrow (+, c 
,M-1),\end{array}\right.\nonumber\\
&&\qquad\qquad\qquad\qquad\quad
\left.\begin{array}{ccc}
\scriptstyle (-, c^{\prime} ,-M) \leftarrow (+, c ,-M),&
\cdots, &
\scriptstyle (-, c^{\prime} ,-1) \leftarrow (+, c ,-1)
\end{array}\right]
\end{eqnarray}
The label $(+, c ,j)$ indicates a downgoing (indicated 
by $+$) $c$-wave (indicated by $c$) with horizontal wavenumber 
$k_{x}=2\pi j/L_{x}$, while  $(-, c^{\prime} ,j)$ 
is a upgoing (indicated 
by $-$) $c^{\prime}$-wave (indicated by $c^{\prime}$) with horizontal wavenumber 
$k_{x}=2\pi j/L_{x}$.
The diagonal elements of (\ref{RdspDef}) are individual plane wave components 
of the downward incident $c$-wave to the upward reflected $c^{\prime}$-wave 
transition.
The array is diagonal, because the medium is laterally invariant on the 
depth interval $z_1 < z < z_2$ so that there is no coupling between the 
various plane wave components. A similar interpretation may be ascribed 
to the other diagonal partitions of   (\ref{Ssparsity}). Of course, when 
the depth interval $z_1 < z < z_2$ is laterally heterogeneous, due to 
{\em plane wave coupling\/}, each of 
the indicated partitions in (\ref{Ssparsity}) is generally fully populated.
An explicit construction for the plane interface reflection/transmission 
arrays is detailed in {\sc \bf  Box} $ {\bf  5.1}$.

%%%%%%%%%%%%%%%%%%%%%%%%%%%%%%%%%%%%%%%%%%%%%%%%%%%%%%%%%%%%%%%%%%%%%%%%%%%%%%%%%%%%%%%%%%%%
%%%%%%%%%%%%%%%%%%%%%%%%%%%%%%%%%%%%%%%%%%%%%%%%%%%%%%%%%%%%%%%%%%%%%%%%%%%%%%%%%%%%%%%%%%%%
%%%%%%%%%%%%%%%%%%%%%%%%%%%%%%%%%%%%%%%%%%%%%%%%%%%%%%%%%%%%%%%%%%%%%%%%%%%%%%%%%%%%%%%%%%%%
\begin{figure}
\fbox{
\begin{minipage} {140.0mm}
{\sc \bf  Box} $ {\bf  5.1}$
{\sc The the plane interface reflection/ transmission arrays} \newline

Let us partition the $6 N_k \times 6 N_k$ array ${\bmi D}_{z}({\bf k}_{z},k_{y})$
from (\ref{Dsuperdef}) into the  $6 N_k \times 3 N_k$ array
\begin{equation} \label{BdDef}
{\bmi B}_{ \downarrow }({\bf 
k}_{x},{k}_{y})=
\left[\begin{array}{ccc}
\left|+,P, {\bf k}_{x},k_{y}  \right\rangle&\left|+,S, {\bf k}_{x},k_{y}  
\right\rangle&\left|+,H, {\bf k}_{x},k_{y}  \right\rangle
\end{array}\right]
\end{equation}
of downgoing partial wave states, and the $6 N_k \times 3 N_k$ array
\begin{equation} \label{BuDef}
{\bmi B}_{ \uparrow }({\bf 
k}_{x},{k}_{y})=
\left[\begin{array}{ccc}
\left|-,P, {\bf k}_{x},k_{y}  \right\rangle&\left|-,S, {\bf k}_{x},k_{y}  
\right\rangle&\left|-,H, {\bf k}_{x},k_{y}  \right\rangle
\end{array}\right]
\end{equation}
of upgoing partial wave states. The component states in (\ref{BdDef}) and 
(\ref{BuDef}) were defined in 
eqns.(\ref{PeigenSuper})--(\ref{SHeigenSuper}).


Let us assume that we have a plane interface problem, where the upper 
and lower half spaces have uniform material properties 
$(\rho^{1}, \alpha^{1}, \beta^{1})$ and 
$(\rho^{2}, \alpha^{2}, \beta^{2})$, respectively. We append 
the superscripts $1$ and $2$ to the arrays in (\ref{BdDef}) and 
(\ref{BuDef}) to indicate that $(\rho, \alpha, \beta)$ in the
definitions (\ref{PeigenSuper})--(\ref{SHeigenSuper}) are to be replaced
by $(\rho^{1}, \alpha^{1}, \beta^{1})$ and 
$(\rho^{2}, \alpha^{2}, \beta^{2})$, respectively. 

The method we employ to construct the plane interface reflection/ 
transmission arrays is straightforward generalization of {\em Kennett\/}~ 
(1983), \S{\bf 5.1}. The canonical displacement/traction balance
\begin{eqnarray} \label{rtbalance}
\lefteqn{\left[{\begin{array}{cc}{\bmi B}_{ \downarrow 
}^{1}({\bf k}_{x},{k}_{y})&{\bmi B}_{ \uparrow 
}^{1}({\bf k}_{x},{k}_{y})\end{array}}\right]
\left[{\begin{array}{cc}{\bmi I}_{3{N}_{k}}&{\bf 0}_{3{N}_{k}}\\ 
{\bmi S}^{1}_{ \uparrow \downarrow }&{\bmi S}^{1}_{ \uparrow \uparrow 
}\end{array}}\right]=\left[{\begin{array}{cc}{\bmi B}_{ \downarrow 
}^{2}({\bf k}_{x},{k}_{y})&{\bmi B}_{ \uparrow 
}^{2}({\bf k}_{x},{k}_{y})\end{array}}\right]\left[{\begin{array}{cc}{\bmi S}^{1}_{ 
\downarrow \downarrow }&{\bmi S}^{1}_{ \downarrow \uparrow }\\ 
{\bf 0}_{3{N}_{k}}&{\bmi I}_{3{N}_{k}}\end{array}}\right]}\nonumber\\
&&
\end{eqnarray} 
expresses the continuity of the $3-$ components of displacement 
$(u_z,u_x,u_y)$ and  
$3-$ components of traction $(\tau_{zz},\tau_{xz},\tau_{yz})$ across the discontinuity of the material 
parameters at the plane interface. We have $3 N_k$ unit amplitude 
downgoing planewaves states incident on the interface from the upper half 
space, and a further  $3 N_k$ unit amplitude 
upgoing planewaves states incident on the interface from the lower half 
space. The $3 N_k \times 3 N_k$ arrays ${\bmi S}^{1}_{ \uparrow \downarrow 
}$, ${\bmi S}^{1}_{ \downarrow \downarrow }$, ${\bmi S}^{1}_{ \uparrow \uparrow 
}$ and ${\bmi S}^{1}_{ \downarrow \uparrow }$ are the undetermined reflection 
and transmission arrays, whose entries are to be adjusted to balance the 
displacement/traction field on the two sides of the boundary.

By rearranging (\ref{rtbalance}) so that the unknown 
reflection and transmission arrays are on the left side of the rearranged 
equations, we get
\begin{equation} \label{rtbalance1}
\left[{\begin{array}{cc}{\bmi B}_{ \uparrow 
}^{1}({\bf k}_{x},{k}_{y})&-{\bmi B}_{ \downarrow 
}^{2}({\bf k}_{x},{k}_{y})\end{array}}\right]\left[{\begin{array}{cc}{\bmi S}^{1}_{ 
\uparrow \downarrow }&{\bmi S}^{1}_{ \uparrow \uparrow }\\ {\bmi S}^{1}_{ \downarrow 
\downarrow }&{\bmi S}^{1}_{ \downarrow \uparrow 
}\end{array}}\right]=\left[{\begin{array}{cc}-{\bmi B}_{ \downarrow 
}^{1}({\bf k}_{x},{k}_{y})&{\bmi B}_{ \uparrow 
}^{2}({\bf k}_{x},{k}_{y})\end{array}}\right]
\end{equation} 

If we introduce the bilinear form $\langle \cdot, \cdot \rangle$ defined by
\[
\langle {\bmi a}, {\bmi b} \rangle= {\bmi a}^{T}\left({\bmi J}_6 \otimes 
{\bmi I}_{N_k}  \right) {\bmi b}
\]
for arbitrary $6 N_k-$vectors ${\bmi a}$ and  ${\bmi b}$ then
\begin{eqnarray} \label{rtbalance2}
\lefteqn{\left[{\begin{array}{c}\begin{array}{cc}{\bf 
0}_{3{N}_{k}}&-\left\langle{{\bmi B}_{ \uparrow }^{1}(-{\bf 
k}_{x},-{k}_{y}),{\bmi B}_{ \downarrow }^{2}({\bf 
k}_{x},{k}_{y})}\right\rangle\end{array}\\ 
\begin{array}{cc}\left\langle{{\bmi B}_{ \downarrow }^{2}(-{\bf 
k}_{x},-{k}_{y}),{\bmi B}_{ \uparrow }^{1}({\bf 
k}_{x},{k}_{y})}\right\rangle&{\bf 
0}_{3{N}_{k}}\end{array}\end{array}}\right]\left[{\begin{array}{cc}{\bmi 
S}^{1}_{ \uparrow \downarrow }&{\bmi S}^{1}_{ \uparrow \uparrow }\\ {\bmi S}^{1}_{ 
\downarrow \downarrow }&{\bmi S}^{1}_{ \downarrow \uparrow 
}\end{array}}\right]=}\nonumber\\
&&
\left[{\begin{array}{c}\begin{array}{cc}-i\,\frac{{\mit \Delta}k_x}{2 
\pi}\,{\bmi 
I}_{3{N}_{k}}&\left\langle{{\bmi B}_{ \uparrow }^{1}(-{\bf 
k}_{x},-{k}_{y}),{\bmi B}_{ \uparrow }^{2}({\bf 
k}_{x},{k}_{y})}\right\rangle\end{array}\\ 
\begin{array}{cc}-\left\langle{{\bmi B}_{ \downarrow }^{2}(-{\bf 
k}_{x},-{k}_{y}),{\bmi B}_{ \downarrow }^{1}({\bf 
k}_{x},{k}_{y})}\right\rangle&-i\,\frac{{\mit \Delta}k_x}{2 \pi}\,{\bmi 
I}_{3{N}_{k}}\end{array}\end{array}}\right] 
\end{eqnarray} 
where we have used
\[
\left\langle{{\bmi B}_{ \uparrow }^{1}(-{\bf k}_{x},-{k}_{y}),{\bmi B}_{ \uparrow 
}^{1}({\bf k}_{x},{k}_{y})}\right\rangle=\left\langle{{\bmi B}_{ \downarrow 
}^{2}(-{\bf k}_{x},-{k}_{y}),{\bmi B}_{ \downarrow 
}^{2}({\bf k}_{x},{k}_{y})}\right\rangle={\bf 0}_{3{N}_{K}}
\] 
to get the left side of (\ref{rtbalance2}) and

\vspace{6mm}
\end{minipage}
}
\end{figure}
\begin{figure}
\fbox{
\begin{minipage} {140.0mm}
{\sc \bf  Box} ${\bf  5.1}$ {\em continued}\newline
\[
\left\langle{{\bmi B}_{ \uparrow }^{1}(-{\bf k}_{x},-{k}_{y}),{\bmi B}_{ 
\downarrow }^{1}({\bf k}_{x},{k}_{y})}\right\rangle=i\, \frac{{\mit \Delta}k_x}{2 \pi}{\bmi I}_{3{N}_{K}}
\] 

along with 
\[
\left\langle{{\bmi B}_{ \downarrow 
}^{2}(-{\bf k}_{x},-{k}_{y}),
{\bmi B}_{ 
\uparrow }^{2}({\bf k}_{x},{k}_{y})}\right\rangle
=-i\, \frac{{\mit \Delta}k_x}{2 \pi}{\bmi I}_{3{N}_{K}}
\] 
to get the right side of (\ref{rtbalance2}).

Finally, by multiplying both sides of  (\ref{rtbalance2}) by
\[
\left[{\begin{array}{c}\begin{array}{cc}{\bf 
0}_{3{N}_{k}}&-\left\langle{{\bmi B}_{ \uparrow }^{1}(-{\bf 
k}_{x},-{k}_{y}),{\bmi B}_{ \downarrow }^{2}({\bf 
k}_{x},{k}_{y})}\right\rangle\end{array}\\ 
\begin{array}{cc}\left\langle{{\bmi B}_{ \downarrow }^{2}(-{\bf 
k}_{x},-{k}_{y}),{\bmi B}_{ \uparrow }^{1}({\bf 
k}_{x},{k}_{y})}\right\rangle&{\bf 
0}_{3{N}_{k}}\end{array}\end{array}}\right]^{-1}\!\!\!=
\left[{\begin{array}{c}\begin{array}{cc}{\bf 
0}_{3{N}_{k}}&{\left\langle{{\bmi B}_{ \downarrow }^{2}(-{\bf 
k}_{x},-{k}_{y}),{\bmi B}_{ \uparrow }^{1}({\bf 
k}_{x},{k}_{y})}\right\rangle}^{-1}\end{array}\\ 
\begin{array}{cc}-{\left\langle{{\bmi B}_{ \uparrow }^{1}(-{\bf 
k}_{x},-{k}_{y}),{\bmi B}_{ \downarrow }^{2}({\bf 
k}_{x},{k}_{y})}\right\rangle}^{-1}&{\bf 
0}_{3{N}_{k}}\end{array}\end{array}}\!\!\right]
\]
we get
\begin{eqnarray} \label{RTdef}
\lefteqn{\left[{\begin{array}{cc}
{\bmi S}^{1}_{ \uparrow \downarrow }&{\bmi S}^{1}_{ 
\uparrow \uparrow }\\ {\bmi S}^{1}_{ \downarrow \downarrow }&{\bmi S}^{1}_{ 
\downarrow \uparrow 
}\end{array}}\right]=\left[{\begin{array}{c}\begin{array}{cc}{\bf 
0}_{3{N}_{k}}&{\left\langle{{\bmi B}_{ \downarrow }^{2}(-{\bf 
k}_{x},-{k}_{y}),{\bmi B}_{ \uparrow }^{1}({\bf 
k}_{x},{k}_{y})}\right\rangle}^{-1}\end{array}\\ 
\begin{array}{cc}-{\left\langle{{\bmi B}_{ \uparrow }^{1}(-{\bf 
k}_{x},-{k}_{y}),{\bmi B}_{ \downarrow }^{2}({\bf 
k}_{x},{k}_{y})}\right\rangle}^{-1}&{\bf 
0}_{3{N}_{k}}\end{array}\end{array}}\right]
}\nonumber\\
&&\qquad\qquad\qquad\qquad \times \left[{\begin{array}{c}
\begin{array}{cc}-i\,\frac{{\mit \Delta}k_x}{2 
\pi}\,{\bmi 
I}_{3{N}_{k}}&\left\langle{{\bmi B}_{ \uparrow }^{1}(-{\bf 
k}_{x},-{k}_{y}),{\bmi B}_{ \uparrow }^{2}({\bf 
k}_{x},{k}_{y})}\right\rangle\end{array}\\ 
\begin{array}{cc}-\left\langle{{\bmi B}_{ \downarrow }^{2}(-{\bf 
k}_{x},-{k}_{y}),{\bmi B}_{ \downarrow }^{1}({\bf 
k}_{x},{k}_{y})}\right\rangle&-i\,\frac{{\mit \Delta}k_x}{2 \pi}\,{\bmi 
I}_{3{N}_{k}}\end{array}\end{array}}\right] 
\end{eqnarray} 

On a serial machine, (\ref{RTdef}) requires the evaluation of a large 
number of $3 \times 3$ matrix inversions. However, on a vector machine 
it is preferable to evaluate (\ref{RTdef}) in terms of its $N_k$-element diagonal 
partitions, because  this approach is highly vectorizable.
\vspace{6mm}
\end{minipage}
}
\end{figure}

%%%%%%%%%%%%%%%%%%%%%%%%%%%%%%%%%%%%%%%%%%%%%%%%%%%%%%%%%%%%%%%%%%%%%%%%%%%%%%%%
%%%%%%%%%%%%%%%%%%%%%%%%%%%%%%%%%%%%%%%%%%%%%%%%%%%%%%%%%%%%%%%%%%%%%%%%%%%%%%%%
%%%%%%%%%%%%%%%%%%%%%%%%%%%%%%%%%%%%%%%%%%%%%%%%%%%%%%%%%%%%%%%%%%%%%%%%%%%%%%%%
%%%%%%%%%%%%%%%%%%%%%%%%%%%%%%%%%%%%%%%%%%%%%%%%%%%%%%%%%%%%%%%%%%%%%%%%%%%%%%%%




\begin{figure}
\HideDisplacementBoxes
%\ShowDisplacementBoxes
\centerline{\BoxedEPSF{Sfig.EPSF scaled 1000}}
\caption{The reflection/transmission operator notation for regions and 
interfaces}
\protect\label{Sfig}
\end{figure}

There are  two common exceptions to this 
nomenclature:
\begin{itemize}
\item for the interface between layers, (say interface-$i$) we write
${\bmi S}_{\uparrow\downarrow }^{i}\equiv {\bmi S}_{\uparrow\downarrow }^{(i-,i+)}$
({\em see\/} Fig.\ref{Sfig}) 
The appended $\pm$ in $i\pm$ indicates
\[
i \pm \equiv  
\lim\limits_{\rm \epsilon\rightarrow 0+} {z}_{i} \pm \epsilon 
\] 

\item for the open  depth interval $z_{i-1} < z < z_{i}$  spanning a single 
layer, we write 
\[{\bmi S}_{\uparrow\uparrow }^{(i-1,i)}=
{\bmi S}_{\downarrow\downarrow }^{(i-1,i)}={\bmi E}_{i}
\]
The $3 N_k\times 3 N_k$ diagonal, plane-layer phase shift array is defined as
\begin{equation} \label{Ejdef}
{\bmi  E}_{i}=\mbox{diag} \left[{\rm e}^{i {\bf  k}_{z, 
P_{i} }(z_{i}-z_{i-1})},\  {\rm e}^{i 
{\bf  k}_{z, S_{i} }(z_{i}-z_{i-1})}
,\  {\rm e}^{i 
{\bf  k}_{z, H_{i} }(z_{i}-z_{i-1})}\right] ,
\end{equation} 
with the $N_k \times N_k$ vertical wavenumber array ${\bf  k}_{z, c_{i} }$ given by
\begin{eqnarray} \label{kzidef} 
{\bf  k}_{z, c_{i} }&\equiv& \mbox{diag}\left[ {k}_{z, c_{i} }(0 ),
{k}_{z, c_{i} }({\mit 
\Delta}k_x),\ \cdots,\right.\nonumber\\
&&\qquad\left. \ {k}_{z, c_{i} }((M-1) {\mit \Delta}k_x)
, \ {k}_{z, c_{i} }(-M {\mit \Delta}k_x)
,\ \cdots, \ {k}_{z, c_{i} }(- {\mit \Delta}k_x) \right]
\end{eqnarray}
for $N_k=2 M$ and $c_i\in \{P_i,S_i, H_i \}$,
where 
\begin{equation}	\label{kzc1Def}
	{k}_{z,c_i}({k}_{x}) = \left\{
	\begin{array}{c}
	\sqrt {\hat{ K}_{c_i}^{2}\ -\ {k}_{x}^{2}}\ \ \ \ \ \ |{k}_{x}|\ <\ 
	|\hat{K}_{c_i}|\\
	i\sqrt {{k}_{x}^{ 2}\ -\ \hat{ K}_{c_i}^{2}} \ \ \ \ |{k}_{x}|\ \ge   \ 
	|\hat{K}_{c_i}|
\end{array}
\right. , \quad  c_i\in\{P_i,S_i,H_i\}, 
\end{equation}
and
\begin{equation}	\label{K1def}
	\hat{K}_{c_i} = \left\{
	\begin{array}{c}
	\sqrt {(\omega_+/{c_i})^{2}\ -\ {k}_{y}^{2}}\ \ \ \ \ \ |{k}_{y}|\ <\ 
	|\omega_+/{c_i}|\\
	i\sqrt {{k}_{y}^{ 2}\ -\ (\omega_+/{c_i})^{2}} \ \ \ \ |{k}_{y}|\ \ge   \ 
	|\omega_+/{c_i}|
\end{array} \right. , \quad  c_i\in\{P_i,S_i,H_i\}, 
\end{equation}
an implicit function of $k_y$.
\end{itemize}


The essence of the {\em standard  multiple scattering formalism\/} [e.g. 
{\em Kennett}~(1986)] is to assume 
that {\em all propagation takes place within a reference medium, with 
the laterally invariant properties\/} ${\bmi  m}_{i}$. The 
material variations ${\mit \Delta}{\bmi  m}_{i}$ of the superimposed heterogeneity 
then induce scattering of the reference medium wave field. There is an 
important variation on the standard theory [e.g. {\em Wu\/}~(1994)]
 that incorporates the forward 
transmission through the reference model and the forward scattered field
from the superimposed heterogeneity directly within in the reference solution, while the 
reflected phases in the reference model and the back scattered field from 
the heterogeneity are governed by a modified Lippmann-Schwinger equation. 
This approach is best suited to problems where forward scattering and 
transmission dominates back scattering and reflection [e.g. for smooth 
reference models at high frequencies]. We concentrate on the standard 
multiple scattering model in this thesis, with details of the forward 
scattering model restricted to \S\ref{fscat}. 

\section[The multiple scattering equations]
{The multiple scattering equations for a stratified halfspace} \label{RTfactor}

Following the  {\em numerical analysis route\/} of {\em Chin, Hedstrom \& 
Thigpen\/}~(1984), propagation through 
the plane stratified component of Fig.\ref{BasicStack}, may be modeled 
via the solution of a linear system of equations that results from the 
enforcement of {\em displacement-traction\/} continuity across the interfacial
boundaries between each plane layer. However, it is not clear how the 
superimposed heterogeneity would be included into the analysis, using the 
method described in {\em Chin et al.\/}~(1984), as the influence of the 
heterogeneity extends over space, and is not confined to interfacial 
boundaries. Instead, we will adopt a {\em multiple scattering\/} 
formalism, where the wavefield is governed by a linear algebraic system that results from 
a set of {\em wavefield balances\/}, rather then the traditional  {\em displacement-traction\/}
balances. 


%%%%%%%%%%%%%%%%%%%%%%%%%%%%%%%%%%%%%%%%%%%%%%%%%%%%%%%%%%%%%%%%%%%%%%%%%%%%%%%%
%%%%%%%%%%%%%%%%%%%%%%%%%%%%%%%%%%%%%%%%%%%%%%%%%%%%%%%%%%%%%%%%%%%%%%%%%%%%%%%%
%%%%%%%%%%%%%%%%%%%%%%%%%%%%%%%%%%%%%%%%%%%%%%%%%%%%%%%%%%%%%%%%%%%%%%%%%%%%%%%%

With reference to Fig.\ref{BasicStack}, the standard multiple scattering model 
is composed from two distinct physical processes:
\begin{enumerate}
\item {\em interfacial wavefield balances\/} (Fig.\ref{IntScat}) 
\begin{equation} \label{ibalance}
\left\{ \begin{array}{ll}
{{\bmi V}}_{ \downarrow }^{1}=
{\bmi S}_{ \downarrow \uparrow}^{[f,{0}]}
{\bmi  u}_{1}&\begin{array}{l}
\mbox{upward incident reflection} \\
\mbox{from the free surface zone}\end{array} \\
\left({\begin{array}{c}
{\bmi  V}_{ \uparrow }^{i}\\
{\bmi  V}_{ \downarrow }^{i}\end{array}}\right)=
\left({\begin{array}{cc}
{\bmi S}_{\uparrow \downarrow }^{i}&{\bmi S}_{ \uparrow \uparrow }^{i}\\ 
{\bmi S}_{\downarrow \downarrow }^{i}&{\bmi S}_{ \downarrow \uparrow}^{i}
\end{array}}\right)
\left({\begin{array}{c}{\bmi  d}_{i}\\
{\bmi  u}_{i}\end{array}}\right)
&1\le i \le L-1\\
{{\bmi V}}_{ \uparrow }^{L}=
{\bmi S}_{ \downarrow \uparrow}^{[L,\infty)}
{\bmi  d}_{L}& \begin{array}{l}
\mbox{downward incident reflection}\\
\mbox{from the lower halfspace}
\end{array}
\end{array} \right.
\end{equation}
to model the transition of the wavefield between adjacent plane layers of the reference 
medium. 
An explicit 
construction for the interfacial reflection/transmission 
arrays is given in 
{\sc \bf  Box} ${\bf  5.1}$.

\begin{samepage}
\item {\em layer wavefield balances\/} (Fig.\ref{LayerScat})
\begin{equation} \label{Lbalance}
\left({\begin{array}{c}{\bmi  u}_{i}\\
{\bmi  d}_{i}\end{array}}\right)=\left({\begin{array}{cc}
{{\mit \Delta} \bmi  S}_{ \uparrow \downarrow}^{i}&{\bmi  E}_{i}+{{\mit \Delta} \bmi  S}_{ \uparrow
\uparrow }^{i}\\ 
{\bmi  E}_{i}+{{\mit \Delta} \bmi  S}_{ \downarrow\downarrow }^{i}&{{\mit \Delta} \bmi  S}_{
\downarrow\uparrow}^{i}
\end{array}}\right)
\left({\begin{array}{c}{\bmi  V}_{ \downarrow }^{i}\\
{\bmi  V}_{ \uparrow }^{i}\end{array}}\right)
\end{equation}
\end{samepage}
to model the plane layer phase shift due to propagation across the 
homogeneous reference layer, as well as the scattered field due to the 
superimposed heterogeneity.
An explicit 
construction for the layer scattering  
${{\mit \Delta} \bmi  S}^{i}$ arrays is given in 
{\sc   Box} ${\bf  5.2}$.

\end{enumerate}

\begin{figure}
\HideDisplacementBoxes
%\ShowDisplacementBoxes
\centerline{\BoxedEPSF{IntScat.EPSF scaled 1000}}
\caption{The wavefield balance for plane interface scattering.}
\protect\label{IntScat}
\end{figure}
\begin{figure}
\HideDisplacementBoxes
%\ShowDisplacementBoxes
\centerline{\BoxedEPSF{LayerScat.EPSF scaled 1000}}
\caption{The wavefield balance for a single heterogeneous layer superimposed on a homogeneous
background medium. The balance involves both background propagation, 
effected by the phase shift $\mbox{\mib  E}{}_{i}$, and the scattering processes 
modeled by ${\mit \Delta}\mbox{\mib S}\,{}^{i}$. }
\protect\label{LayerScat}
\end{figure}

%%%%%%%%%%%%%%%%%%%%%%%%%%%%%%%%%%%%%%%%%%%%%%%%%%%%%%%%%%%%%%%%%%%%%%%%%%%%%%%%
%%%%%%%%%%%%%%%%%%%%%%%%%%%%%%%%%%%%%%%%%%%%%%%%%%%%%%%%%%%%%%%%%%%%%%%%%%%%%%%%
%%%%%%%%%%%%%%%%%%%%%%%%%%%%%%%%%%%%%%%%%%%%%%%%%%%%%%%%%%%%%%%%%%%%%%%%%%%%%%%%
%%%%%%%%%%%%%%%%%%%%%%%%%%%%%%%%%%%%%%%%%%%%%%%%%%%%%%%%%%%%%%%%%%%%%%%%%%%%%%%%
\begin{figure}
\fbox{
\begin{minipage} {140.0mm}
{\sc   Box} ${\bf  5.2}$
{\sc The layerstack scattering matrix} \newline

We define the sideways propagator 
\begin{eqnarray}\label{CBPxaug}
\!\!\!\!\!\!\!\!\!
\lefteqn{\rf{\CB P}\!{}^{x}_{r,s,\rho} =}\nonumber\\
&&\!\!\!\!\!\!\!\!\!
 -{i} {\bmi I}_{N_x} \otimes
\left\{\left({\bmi I}_{8}\otimes \overline{\bmi F}(0)  \right)
 \hat{\bmi  D}_{x,\rho}({\bf k}_{z},{k}_{y})\right\}
 {\CB Q}{}^{x}_{\rho \leftarrow \rho}{\bmi I}_{N_x} \otimes 
 \left\{\hat{\bmi  D}_{x,\rho}^{T}(-{\bf k}_{z},-{k}_{y})
\,\left({\bmi I}_{8}\otimes \overline{\bmi F}\!{}^{T}(0)  \right)
\right\}
\end{eqnarray}
for $\mbox{row}_{\rho}$ of the heterogeneity lattice,
where the $8 \bar{N}_k \times 6 \bar{N}_k$ array 
$\hat{\bmi  D}_{x,\rho}$  is an augmented version
 of the $6 \bar{N}_k \times 6 \bar{N}_k$ array 
${\bmi  D}_{x}$ (\ref{Dsuperdef}) and is given in {\sc Box } {\bf 4.6}. 
The appended $\rho$ in $\hat{\bmi  D}_{x,\rho}$ indicates that it should 
be constructed using the reference material parameters appropriate to 
$\mbox{row}_{\rho}$.
The planewave propagator ${\CB Q}{}^{x}_{\rho \leftarrow \rho}$
was given by (\ref{CBQxdef}).
The projection matrix $\overline{\bmi F}(0)$ was given by 
(\ref{Fbardef}).

$\quad$Using (\ref{CBPxaug}) we can sum all possible sideways multiple scattering 
interactions between the $N_x$ squares of heterogeneity superimposed
on  layer-$\rho$,  by introducing 
the $6 N_x \times 6 N_x$ layer transition matrix, defined as
\begin{equation} \label{layerT}
{\mit \Delta}{\CB T}{}^{\rho}(k_y)=
{\mit \Delta}{\CB C}{}^{\rho}_{eff}(k_y)\left[{\bmi I}- \rf{\CB P}\!{}^{x}_{r,s,\rho}\cdot {\mit \Delta}{\CB 
C}{}^{\rho}_{eff}(k_y)  \right]^{-1}, \quad 1 \le \rho \le L,
\end{equation}
where  the {\em effective coupling 
matrix\/} ${\mit \Delta}{\CB C}{}^{\rho}_{eff}(k_y)$ for $\mbox{row}_{\rho}$
was defined by (\ref{dCBCrho}).

 Using (\ref{layerT}) , we define
the $6 N_k \times 6 N_k$ array of plane wave reflection/transmission 
transitions from $\mbox{row}_{\rho}$ of the heterogeneity by
\begin{eqnarray} \label{laySdef}
\lefteqn{
{\mit \Delta}{\CB S}^{\rho}(k_y)=
\left({\begin{array}{cc}
{\mit \Delta}{\bmi S}_{\uparrow \downarrow }^{\rho}&{\mit \Delta}{\bmi S}_{ \uparrow \uparrow }^{\rho}\\ 
{\mit \Delta}{\bmi S}_{\downarrow \downarrow }^{\rho}&{\mit \Delta}{\bmi S}_{ \downarrow \uparrow}^{\rho}
\end{array}}\right)=
-i\,\left({\begin{array}{cc}
\sqrt{{\bmi E}_{\rho}}&{\bf 0}\\ 
{\bf 0}&\sqrt{{\bmi E}_{\rho}}
\end{array}}\right)\,
\hat{\bmi  D}{}^{T}_{z,\rho}(-{\bf k}_{x},-{k}_{y})}\\
&&
\times 
\left({\bmi I}_8\otimes{\bf DFT}^{T}(-{\bf k}_{x})  \right)
\hat{\bmi \Phi}{\mit \Delta}{\bmi T}{}^{\rho}(k_y) 
\hat{\bmi \Phi}{}^{-1} \left({\bmi I}_8\otimes 
{\bf DFT}({\bf k}_{x}) \right)\hat{\bmi  D}_{z,\rho}({\bf k}_{x},{k}_{y}) \left({\begin{array}{cc}
\sqrt{{\bmi E}_{\rho}}&{\bf 0}\\ 
{\bf 0}&\sqrt{{\bmi E}_{\rho}}
\end{array}}\right)\nonumber
\end{eqnarray}
which incorporates the upsampling $\hat{\bmi \Phi}$ and downsampling $\hat{\bmi \Phi}{}^{-1}$ 
operators that were defined in Alg 4.1.
The appended $\rho$ in $\hat{\bmi  D}_{z,\rho}$ indicates that it should 
be constructed using the reference material parameters appropriate to 
$\mbox{row}_{\rho}$ according to its definition {\sc Box } {\bf 4.5}.
Also, 
\[
\left({\begin{array}{cc}
\sqrt{{\bmi E}_{\rho}}&{\bf 0}\\ 
{\bf 0}&\sqrt{{\bmi E}_{\rho}}
\end{array}}\right)
\]
implements the  {\em phase shift\/} of the interface variables ${\CB  
V}_{{int}}$ of (\ref{Xdef})  from their 
assumed phase reference values at the edges of layer-$\rho$,
 to the centre of that 
layer. We recall that ${\bmi E}_{\rho}$ was defined in (\ref{Ejdef}).



We now assemble the blocks (\ref{laySdef}) for each $1 \le \rho \le L$ 
into the block diagonal {\em layer stack scattering array\/}
 ${\mit \Delta}{\CB  S}(k_y)$ defined as
\begin{equation} \label{GSdef}
{\mit \Delta}{\CB  S}(k_y) \ \equiv\ \left({\begin{array}{ccc}
\begin{array}{cc}{\mit \Delta}{\bmi S}_{\uparrow
\downarrow }^{1}&{\mit \Delta}{\bmi S}_{\uparrow \uparrow }^{1}\\ 
{\mit \Delta}{\bmi S}_{\downarrow\downarrow }^{1}&{\mit \Delta}{\bmi S}_{\downarrow \uparrow }^{1}
\end{array}&&\\ 
& \ddots &\\  
&&\begin{array}{cc}{\mit \Delta}{\bmi S}_{\uparrow \downarrow}^{L}&
{\mit \Delta}{\bmi S}_{\uparrow \uparrow }^{L}\\ 
{\mit \Delta}{\bmi S}_{\downarrow \downarrow }^{L}
&{\mit \Delta}{\bmi S}_{\downarrow \uparrow }^{L}\end{array}
\end{array}}\right)
\end{equation}


 




\vspace{6mm}
\end{minipage}
}
\end{figure}

\begin{figure}
\fbox{
\begin{minipage} {140.0mm}
{\sc   Box}  ${\bf  5.2}$
{\em continued} \newline

For computational purposes, it will be convenient to factor the block 
diagonal array  ${\mit \Delta}{\CB  S}(k_y)$ of (\ref{GSdef}) into the 
product of separate block diagonal arrays. To this end,
we define the block diagonal arrays:
\begin{eqnarray} \label{Thetaouz}
\lefteqn{{\bmi \Theta}^{z}_{ou}\equiv
\mbox{diag}\left[
\left({\begin{array}{cc}
\sqrt{{\bmi E}_{1}}&{\bf 0}\\ 
{\bf 0}&\sqrt{{\bmi E}_{1}}
\end{array}}\right)\,\hat{\bmi  D}{}^{T}_{z,1}(-{\bf k}_{x},-{k}_{y})\left({\bmi I}_8\otimes
{\bf DFT}^{T}(-{\bf k}_{x})  \right)\hat{\bmi \Phi} ,\right.}\nonumber\\
&& \hspace{6em}
\left. \cdots,
\left({\begin{array}{cc}
\sqrt{{\bmi E}_{L}}&{\bf 0}\\ 
{\bf 0}&\sqrt{{\bmi E}_{L}}
\end{array}}\right)\,\hat{\bmi  D}{}^{T}_{z,L}(-{\bf k}_{x},-{k}_{y})\left({\bmi I}_8\otimes
{\bf DFT}^{T}(-{\bf k}_{x})  \right)\hat{\bmi \Phi}   
\right];
\end{eqnarray}

\begin{equation} \label{CBT}
{\mit \Delta}{\CB T}(k_y)\equiv
\mbox{diag}\left[
{\mit \Delta}{\CB T}{}^{1}(k_y), \cdots,
{\mit \Delta}{\CB T}{}^{L}(k_y)
\right];
\end{equation}
\begin{eqnarray} \label{Thetainz}
\lefteqn{{\bmi \Theta}^{z}_{in}\equiv
\mbox{diag}\left[
\hat{\bmi \Phi}{}^{-1}\left({\bmi I}_8\otimes 
{\bf DFT}({\bf k}_{x}) \right)\hat{\bmi  D}_{z,1}({\bf k}_{x},{k}_{y})
\left({\begin{array}{cc}
\sqrt{{\bmi E}_{1}}&{\bf 0}\\ 
{\bf 0}&\sqrt{{\bmi E}_{1}}
\end{array}}\right),\right.}\nonumber\\
&&\hspace{9em}
\left. \cdots,
\hat{\bmi \Phi}{}^{-1}\left({\bmi I}_8\otimes 
{\bf DFT}({\bf k}_{x}) \right)\hat{\bmi  D}_{z,L}({\bf k}_{x},{k}_{y})
\left({\begin{array}{cc}
\sqrt{{\bmi E}_{L}}&{\bf 0}\\ 
{\bf 0}&\sqrt{{\bmi E}_{L}}
\end{array}}\right)
\right].
\end{eqnarray}


Using (\ref{Thetaouz})--(\ref{Thetainz}),
 the layer stack scattering array ${\mit \Delta} {\CB  S}$ of 
 (\ref{GSdef}) may be 
written  as
\begin{equation} \label{GSdef1}
{\mit \Delta}{\CB  S}(k_y) \ \equiv\ \left({\begin{array}{ccc}
\begin{array}{cc}{\mit \Delta}{\bmi S}_{\uparrow
\downarrow }^{1}&{\mit \Delta}{\bmi S}_{\uparrow \uparrow }^{1}\\ 
{\mit \Delta}{\bmi S}_{\downarrow\downarrow }^{1}&{\mit \Delta}{\bmi S}_{\downarrow \uparrow }^{1}
\end{array}&&\\ 
& \ddots &\\  
&&\begin{array}{cc}{\mit \Delta}{\bmi S}_{\uparrow \downarrow}^{L}&
{\mit \Delta}{\bmi S}_{\uparrow \uparrow }^{L}\\ 
{\mit \Delta}{\bmi S}_{\downarrow \downarrow }^{L}
&{\mit \Delta}{\bmi S}_{\downarrow \uparrow }^{L}\end{array}
\end{array}}\right)
\equiv -i\, {\bmi \Theta}^{z}_{ou} {\mit \Delta}{\CB T}(k_y) {\bmi \Theta}^{z}_{in}
\end{equation}

\vspace{6mm}
\end{minipage}
}
\end{figure}
%%%%%%%%%%%%%%%%%%%%%%%%%%%%%%%%%%%%%%%%%%%%%%%%%%%%%%%%%%%%%%%%%%%%%%%%%%%%%%%%
%%%%%%%%%%%%%%%%%%%%%%%%%%%%%%%%%%%%%%%%%%%%%%%%%%%%%%%%%%%%%%%%%%%%%%%%%%%%%%%%
%%%%%%%%%%%%%%%%%%%%%%%%%%%%%%%%%%%%%%%%%%%%%%%%%%%%%%%%%%%%%%%%%%%%%%%%%%%%%%%%
%%%%%%%%%%%%%%%%%%%%%%%%%%%%%%%%%%%%%%%%%%%%%%%%%%%%%%%%%%%%%%%%%%%%%%%%%%%%%%%%
%%%%%%%%%%%%%%%%%%%%%%%%%%%%%%%%%%%%%%%%%%%%%%%%%%%%%%%%%%%%%%%%%%%%%%%%%%%%%%%%
%%%%%%%%%%%%%%%%%%%%%%%%%%%%%%%%%%%%%%%%%%%%%%%%%%%%%%%%%%%%%%%%%%%%%%%%%%%%%%%%

In particular, we note that
${ \bmi  S}^{i}$ effects the reflection and transmission transitions of 
the field for propagation between adjacent layers of the reference 
stratification, while 
 ${{\mit \Delta} 
\bmi  S}^{i}$ effects the additional shift in phase and distortion of amplitude
due to scattering from the superimposed heterogeneities, 
{\em above and beyond\/} the phase shift (modeled by ${\bmi E}_i$) for 
propagation across the homogeneous background. For this reason, ${{\mit \Delta} 
\bmi  S}^{i}$ is referred to as a {\em transition matrix\/} in formal 
scattering theory, and not as a {\em scattering matrix\/}, although the 
term {\em residual scattering matrix\/} might be more informative, as it 
models the scattering from the material variations that are left over, 
after we have removed the effect of the background model. 
The ${\mit \Delta}$ is used to remind us, that ${\mit \Delta}{\bmi S}^i\rightarrow{\bf 0}$ 
as ${\mit \Delta}{\bmi  m}_i\rightarrow{\bf 0}$.


Let us assume the following ordering of the field variables $\CB V$:
\begin{equation} \label{Xorder}
{\CB  V}=\left({\begin{array}{c}{\CB  V}_{{int}}\\
{\CB  V}_{{layer}}\end{array}}\right)
\end{equation}
where (Fig.\ref{BasicStack})
\begin{equation} \label{Xdef}
{\CB  V}_{{int}}=
\left({\begin{array}{c}
{{\bmi V}}_{ \downarrow }^{1}\\
{{\bmi V}}_{ \uparrow }^{1}\\
{{\bmi V}}_{ \downarrow }^{2}\\
\vdots\\
{{\bmi V}}_{ \uparrow }^{L}
\end{array}}\right), \qquad
{\CB  V}_{{layer}}=
\left({\begin{array}{c}
{\bmi  u}_{1}\\
{\bmi  d}_{1}\\
{\bmi  u}_{2}\\
\vdots\\
{\bmi  d}_{L}
\end{array}}\right).
\end{equation}


When the balances (\ref{ibalance}) and (\ref{Lbalance}) are  assembled into a matrix 
consistent with the adopted ordering (\ref{Xorder}) of the field variables,
we get
\begin{equation} \label{Gbalance}
\left({\begin{array}{cc}
{\bmi  I}_{6 N_k L}&-{\CB  S}_{{int}}\\ 
-{\CB  S}_{{layer}}&
{\bmi  I}_{6 N_k L}\end{array}}\right)\left({\begin{array}{c}{\CB  V}_{{int}}\\ 
{\CB  V}_{{layer}}\end{array}}\right)=\left({\begin{array}{c}{\bmi  b}_{{int}}\\ 
{\bf 0}_{6 N_k L \times 1}\end{array}}\right)
\end{equation} 
where the $(6 N_k L) \times (6 N_k L)$ partition ${\CB  S}_{{int}}$ is 
written:
\begin{equation} \label{Sint}
{\CB  S}_{{int}} \ =\ \left({\begin{array}{ccccc}
{\bmi S}_{\downarrow\uparrow}^{[f,0]}&&&&\\ 
&\begin{array}{cc}{\bmi S}_{\uparrow
\downarrow }^{1}&{\bmi S}_{\uparrow \uparrow }^{1}\\ 
{\bmi S}_{\downarrow
\downarrow }^{1}&{\bmi S}_{\downarrow \uparrow 
}^{1}\end{array}&&&\\ && \ddots
&&\\  &&&\begin{array}{cc}{\bmi S}_{\uparrow \downarrow 
}^{L-1}&{\bmi S}_{\uparrow \uparrow }^{L-1}\\ 
{\bmi S}_{\downarrow \downarrow 
}^{L-1}&{\bmi S}_{\downarrow \uparrow 
}^{L-1}\end{array}&\\ 
&&&&{\bmi S}_{\uparrow\downarrow}^{[L,\infty)}\end{array}}\right);
\end{equation}
and  the $(6 N_k L) \times (6 N_k L)$ partition ${\CB  S}_{{layer}}$ is 
defined
\begin{equation} \label{Slayer}
{\CB  S}_{{layer}} \ =\ \left({\begin{array}{ccc}
\begin{array}{cc}{\mit \Delta}{\bmi S}_{\uparrow
\downarrow }^{1}&{\bmi  E}_{1}+{\mit \Delta}{\bmi S}_{\uparrow \uparrow }^{1}\\ 
{\bmi  E}_{1}+{\mit \Delta}{\bmi S}_{\downarrow\downarrow }^{1}&{\mit \Delta}{\bmi S}_{\downarrow \uparrow }^{1}
\end{array}&&\\ 
& \ddots &\\  
&&\begin{array}{cc}{\mit \Delta}{\bmi S}_{\uparrow \downarrow 
}^{L}&{\bmi  E}_{L}+{\mit \Delta}{\bmi S}_{\uparrow \uparrow }^{L}\\ 
{\bmi  E}_{L}+{\mit \Delta}{\bmi S}_{\downarrow \downarrow }^{L}
&{\mit \Delta}{\bmi S}_{\downarrow \uparrow }^{L}\end{array}
\end{array}}\right)
\end{equation}
with ${\bmi  E}_{j}$ given by (\ref{Ejdef}), for
the {\em phase shift\/} across layer-$j$ of the reference medium. 

Finally, the source vector ${\bmi  b}_{{int}}$ in (\ref{Gbalance})  
depends on the location and nature of the source point. For simplicity, 
we restrict ourselves to problems with the source in the free surface 
layer $f\le z \le z_0$ or the lower halfspace $z_L \le z \le \infty$.
For a source in the surface zone (Fig.\ref{sourceS}), we get
\begin{equation} \label{SurfSdef}
{\bmi  b}_{{int}}=\left({\begin{array}{c}{\bmi S}_{ \downarrow \downarrow
}^{(s,0]}{\left({{\bmi  I}-{\bmi S}_{ \downarrow \uparrow }^{[f,s)}{\bmi S}_{
\downarrow\uparrow }^{(s,0]}}\right)}^{-1}\left({{\bmi  \Sigma}_{ \downarrow
}+{\bmi S}_{
\downarrow \uparrow }^{[f,s)}{\bmi  \Sigma}_{ \uparrow }}\right)\\ {\bf 0}\\
\vdots\end{array}}\right)
\end{equation}
and for a source in the lower halfspace (Fig.\ref{sourceL}), we have
\begin{equation} \label{LowerSdef}
{\bmi  b}_{{int}}=
\left({\begin{array}{c}\vdots\\ {\bf 0}\\ {\bmi S}_{ \uparrow \uparrow}^{[L,s)}
{\left({{\bmi  I}-{\bmi S}_{ \uparrow \downarrow }^{(s,\infty)}{\bmi S}_{
\downarrow\uparrow }^{[L,s)}}\right)}^{-1}\left({{\bmi  \Sigma}_{ \uparrow
}+{\bmi S}_{\uparrow \downarrow }^{(s,\infty)}{\bmi  \Sigma}_{ \downarrow}}\right)
\end{array}}\right)
\end{equation}
where the upgoing source radiation is given by ${\bmi \Sigma}_{\uparrow}$, 
and the downgoing source radiation is given by ${\bmi \Sigma}_{\downarrow}$
({\em see\/}{\sc Box} $\bf 5.3$). 

\begin{figure}
%\vspace*{10em}
\fbox{
\begin{minipage} {140.0mm}
{ \sc Box} $\bf 5.3$
{\sc Calculating the plane wave source radiation ${\bmi \Sigma}_{\uparrow}$, 
${\bmi \Sigma}_{\downarrow}$ given a dipolar body
${\bmi F}$}\newline
\vspace{4mm}
We assume the general form
\[
{\bmi F}({k}_{x},{k}_{y};z)={\bmi F}_{1}({k}_{x},{k}_{y};z)\delta(z-{z}_{S})+
{\bmi F}_{2}({k}_{x},{k}_{y};z){\delta }^{\prime}(z-{z}_{S})
\]
for a {\em dipolar point force\/} in the intersection of the $y=0$,  $z=z_S$ 
planes,  which is essentially {\em Kennett\/}~(1983), Eqn.(2.78).\newline
$\quad$Introducing the bilinear form $\langle \cdot , \cdot \rangle$ defined as 
(see also (\ref{functional}))
\[
\langle {\bmi  b}_{1} , {\bmi  b}_{2} \rangle \stackrel{\mbox{\tiny def}}{=}
({\bmi  b}_{1})^{ T} {\bmi  J}_6 {\bmi  b}_{2}
\]
for arbitrary $6-$vector fields ${\bmi  b}_{1}$ and ${\bmi  
b}_{2}$, 
the ${\stackrel{{\mbox{{\small upgoing }}}}{{}_{\mbox{{\small downgoing 
}}}}}$ source radiation $\left[{\begin{array}{c}
{\bmi  \Sigma}_{ \uparrow }({k}_{x},{k}_{y})\\
{\bmi  \Sigma}_{ \downarrow }({k}_{x},{k}_{y})\end{array}}\right]$ at 
fixed horizontal wavenumber $({k}_{x},{k}_{y})$, may be witten as
\begin{equation} \label{source}
\left[{\begin{array}{c}
{\bmi  \Sigma}_{ \uparrow }({k}_{x},{k}_{y})\\
{\bmi  \Sigma}_{ \downarrow }({k}_{x},{k}_{y})\end{array}}\right]=- i\,
\left\langle{{\bmi  D}_{z,S}(-{k}_{x},-{k}_{y}),{\bmi F}_{1}({k}_{x},{k}_{y};z)+
\rf{\bmi  A}_{S}({k}_{x},{k}_{y}){\bmi F}_{2}({k}_{x},{k}_{y};z)}
\right\rangle
\end{equation}
where $x_{S}$ is the horizontal offset of the source point, 
 ${\bmi  D}_{z,S}({k}_{x},{k}_{y})$ and $\rf{\bmi  A}_{S}({k}_{x},{k}_{y})$ 
were given by (\ref{Ddef}) and (\ref{Akdef}) respectively, and the appended ``$S$" 
subscript indicates that they should be constructed from the reference 
material properties $(\rf{\rho}, \rf{\alpha}, \rf{\beta})$ on the source 
plane $z=z_S$.

To demonstrate the application of (\ref{source}) we find the source 
radiation generated by a point explosion at $(z,x,y)=(z_S,x_S,0)$. The 
{\em volume force equivalent\/} of a point explosion in the $(k_x, k_y, \omega)$ 
domain may be written as
\[
{\bmi F}({k}_{x},{k}_{y};z)=M(\omega) \left[{\begin{array}{c}0\\0\\0\\0\\
i {k}_{x}\exp \left({-i {k}_{x}{x}_{S}}\right)\\
i {k}_{y}\exp \left({-i {k}_{x}{x}_{S}}\right)\end{array}}\right]
 \delta(z-{z}_{S})+ M(\omega)
 \left[{\begin{array}{c}0\\0\\0\\\exp \left({-i{k}_{x}{x}_{S}}\right)\\
0\\
0\end{array}}\right]{\delta }^{\prime}(z-{z}_{S})
\]
where $M(\omega)$ is the source spectrum.
Consequently,
\begin{eqnarray*}
{\bmi F}_{1}({k}_{x},{k}_{y};z)+
\rf{\bmi  A}_{S}({k}_{x},{k}_{y}){\bmi 
F}_{2}({k}_{x},{k}_{y};z)&=&M(\omega)\, \exp \left({-i 
{k}_{x}{x}_{S}}\right) {\delta }(z-{z}_{S})\\
&&\times
\left[ \begin{array}{cccccc}{1\over {{\rf{\rho}\,\rf{\alpha}\!{}^2} }}, &  0, 
& 0, & 0, & 
{{2\,i\,{\it k_x}\,{\rf{\beta}\!{}^2}}\over {{\rf{\alpha}\!{}^2}}}, &
  {{2\,i\,{\it k_y}\,{\rf{\beta}\!{}^2}}\over {{\rf{\alpha}\!{}^2}}}
  \end{array}\right]^{T}
\end{eqnarray*}
so that
\[
\left[{\begin{array}{c}
{\bmi  \Sigma}_{ \uparrow }({k}_{x},{k}_{y})\\
{\bmi  \Sigma}_{ \downarrow }({k}_{x},{k}_{y})\end{array}}\right]
={\frac{-i\, \omega\,  M( \omega)\,\exp \left({-i 
{k}_{x}{x}_{S}}\right)}{\rf{\alpha}\!{}^2\,
\sqrt {2\rf{\rho} { k}_{ z, P_S}({k}_{x},{k}_{y})}}}
{\left[{\begin{array}{cccccc}1,&0,&0,&1,&0,&0\end{array}}\right]}^{T}
\] 
where ${ k}_{ z, P_S}({k}_{x},{k}_{y})$ was given by (\ref{kzc1Def}).
\vspace{6 mm}
\end{minipage}
}
\end{figure}


\begin{figure}
\HideDisplacementBoxes
%\ShowDisplacementBoxes
\centerline{\BoxedEPSF{sourceS.EPSF scaled 1000}}
\caption{The incident field for a point source in the surface zone $f\le 
z_s \le z_0$}
\protect\label{sourceS}
\end{figure}

\begin{figure}
\HideDisplacementBoxes
%\ShowDisplacementBoxes
\centerline{\BoxedEPSF{sourceL.EPSF scaled 1000}}
\caption{The incident field for a point source in the lower halfspace
$z_s \ge z_L$}
\protect\label{sourceL}
\end{figure}





Before we attempt to invert the left hand side of (\ref{Gbalance}), we eliminate the 
layer variables ${\CB  V}_{{layer}}$, to get the following {\em Schur complement system\/} 
 for the interface variables ${\CB  V}_{{int}}$
\begin{equation} \label{schur}
({\bmi  I}-{\CB  S}_{{int}}{\CB  S}_{{layer}}){\CB  V}_{{int}}={\bmi b}_{{int}}
\end{equation}
%%%%%%%%%%%%%%%%%%%%%%%%%%%%%%%%%%%%%%%%%%%%%%%%%%%%%%%%%%%%%%%%%%%%%%%%%%%%%%%%%%%%%
The linear system (\ref{schur}) is much too large to be to solved via 
Gaussian based matrix factorization because the diagonal blocks of ${\CB  
S}_{{layer}}$ are fully populated  $(6 N_k) \times (6 N_k)$ complex arrays (for the 
problems tested in Chapter \ref{variational}, $6 N_k \approx 24\times 10^3$). 
In this thesis we propose to {\em split\/} the left hand side array in 
(\ref{schur}) into the sum of two parts. One part is then transferred to 
the right hand side of the equation, while the other part is numerically 
inverted.

\subsection{The standard splitting and the recursive Riccati algorithm}
\label{Riccati} 

We split the {\em layer \/} term ${\CB  S}_{{layer}}$ of 
(\ref{Slayer}) into the sum
\begin{equation} \label{Ssplit}
{\CB  S}_{{layer}}={\CB  E}+ {\mit \Delta}{\CB  S}
\end{equation}
where
\begin{equation} \label{GEdef}
{\CB  E} \ =\ \left({\begin{array}{ccc}
\begin{array}{cc}{\bf 0}&{\bmi  E}_{1}\\ 
{\bmi  E}_{1}&{\bf 0}
\end{array}&&\\ 
& \ddots &\\  
&&\begin{array}{cc}{\bf 0}&{\bmi  E}_{L}\\ 
{\bmi  E}_{L}
&{\bf 0}\end{array}
\end{array}}\right);
\end{equation}
is {\em very sparse\/}, but the indicated partitions of 
${\mit \Delta}{\CB  S}$ defined in (\ref{GSdef}) of {\sc Box} {\bf 5.2} are {\em dense\/}.
We now transfer the ${\mit \Delta}{\CB  S}$ terms to the right hand side of 
(\ref{schur}) to get
\begin{equation} \label{eqn100}
({\bmi  I}-{\CB  S}_{{int}}{\CB  E}){\CB  V}_{{int}}={\bmi  
b}_{{int}}+{\CB  S}_{{int}}{\mit \Delta}{\CB  S}{\CB  V}_{{int}}
\end{equation}
where

\begin{samepage}
\begin{eqnarray} \label{mollref}
\lefteqn{ 
({\bmi  I}-{\CB  S}_{{int}}{\CB  E})=}\\
&&\!\!\!\!\!\!
\left[{\begin{array}{cccccccc}
 {\bmi  I}& \!\!\! -{\bmi  S}_{ \downarrow \uparrow }^{[f,0]}{\bmi  E}_{1}
 & \!\!\!& \!\!\!& \!\!\!& \!\!\!& \!\!\!& \!\!\!\\  
-{\bmi  S}_{\uparrow \downarrow }^{1}{\bmi  E}_{1}& \!\!\!  {\bmi  I}& \!\!\! {\bf 0}& \!\!\!
-{\bmi  S}_{ \uparrow \uparrow }^{1}{\bmi  E}_{2}& \!\!\!& \!\!\!& \!\!\!& \!\!\!\\  
-{\bmi  S}_{ \downarrow \downarrow }^{1}{\bmi  E}_{1}& \!\!\! {\bf 0}& \!\!\!
 {\bmi  I}& \!\!\!-{\bmi  S}_{ \downarrow \uparrow }^{1}
{\bmi  E}_{2}& \!\!\!& \!\!\!& \!\!\!& \!\!\!\\ \!\!\! 
& \!\!\!& \!\!\!\ddots& \!\!\!& \!\!\!& \!\!\!& \!\!\!& \!\!\!\\ \!\!\! 
& \!\!\!& \!\!\!-{\bmi  S}_{ \uparrow \downarrow }^{i}{\bmi  E}_{i}& \!\!\!
 {\bmi  I}& \!\!\!{\bf 0}& \!\!\!-{\bmi  S}_{ \uparrow \uparrow }^{i}{\bmi  E}_{i+1}& \!\!\!& \!\!\!\\ \!\!\! 
& \!\!\!& \!\!\! -{\bmi  S}_{ \downarrow \downarrow}^{i}{\bmi  E}_{i}& \!\!\!{\bf 0}& \!\!\!{\bmi  I}& \!\!\!
-{\bmi  S}_{ \downarrow \uparrow }^{i}{\bmi  E}_{i+1}& \!\!\!& \!\!\!\\ \!\!\! 
& \!\!\!& \!\!\!& \!\!\!& \!\!\! \ddots& \!\!\!& \!\!\!& \!\!\!\\ \!\!\! 
& \!\!\!& \!\!\!& \!\!\!& \!\!\!-{\bmi  S}_{ \uparrow \downarrow }^{L-1}{\bmi  E}_{L-1}& \!\!\!{\bmi 
I}& \!\!\! {\bf 0}& \!\!\! -{\bmi  S}_{ \uparrow \uparrow }^{L-1}{\bmi  E}_{L}\\ \!\!\! 
& \!\!\!& \!\!\!& \!\!\!& \!\!\! -{\bmi  S}_{ \downarrow \downarrow}^{L-1}{\bmi  E}_{L-1}& \!\!\! 
{\bf 0}& \!\!\! {\bmi  I}& \!\!\!
 -{\bmi  S}_{ \downarrow \uparrow }^{L-1}{\bmi  E}_{L}\\ \!\!\! 
& \!\!\!& \!\!\!& \!\!\!& \!\!\!& \!\!\!& \!\!\!  -{\bmi  S}_{ \uparrow 
\downarrow}^{[L,\infty)}{\bmi  E}_{L}& \!\!\!
 {\bmi  I}\end{array}}\right] \nonumber
\end{eqnarray}
\end{samepage}
This splitting is best suited to a problems with weak heterogeneity 
superimposed on a plane stratified  reference model that may change rapidly 
on the scale of the incident wavelength.


The left hand side array in (\ref{eqn100}) governs the wavefield in the 
plane stratified reference model, and it has a sparsity pattern that 
affords an efficient matrix inversion. Consequently, we invert the left hand side 
of (\ref{eqn100}) to get the {\em matrix 
analog of the elastodynamic Lippmann-Schwinger equation\/}
\begin{equation} \label{LSmat}
\fbox{$\displaystyle
{\CB  V}_{{int}}=({\bmi  I}-{\CB  S}_{{int}}{\CB  E})^{-1}{\bmi  b}_{{int}}+
({\bmi  I}-{\CB  S}_{{int}}{\CB  E})^{-1}{\CB  S}_{{int}} {\mit \Delta} 
{\CB  S}{\CB  V}_{{int}}
$}
\end{equation}
or
\[
{\CB  V}_{{int}}
={\CB  V}_{{inc}} + {\CB Q}^{\partial} {\mit \Delta} {\CB  S}{\CB  V}_{{int}}
\]
where
\begin{equation}	\label{PstratDef}
	{\CB Q}^{\partial}\equiv 
	({\bmi  I}-{\CB  S}_{{int}}{\CB  E})^{-1}{\CB  S}_{{int}}
\end{equation}
is the {\em required extension of the plane wave propagator  (\ref{CBQdef}) to a plane stratified 
medium\/}, and
\begin{equation}	\label{sourceV}
	{\CB  V}_{{inc}}\equiv 
	({\bmi  I}-{\CB  S}_{{int}}{\CB  E})^{-1}{\bmi  b}_{{int}}
\end{equation}
is the {\em incident field\/} propagating through the reference medium 
from the source point.The  explicit form of ${\CB  V}_{{inc}}$
 depends on the source vector ${\bmi b}_{{int}}$ of (\ref{SurfSdef}) 
 or (\ref{LowerSdef}) depending on whether the source point is in the surface 
zone, or the lower halfspace respectively.
 For a source in the surface zone, the incident field in layer-$i$ is 
 given by 
\begin{eqnarray} \label{incdown}
\lefteqn{
{\bmi  V}_{{inc}}^{ i}=\left[{\begin{array}{c}{\bmi  V}_{{inc},\downarrow}^{ 
i}\\ 
{\bmi  V}_{{inc},\uparrow}^{ i}\end{array}}\right]=
\left[{\begin{array}{c}{\bmi  I}_{ 3 N_k}\\ 
{\bmi S}{}_{\uparrow\downarrow }^{[i, 
\infty)}{{\bmi E}}_{i}\end{array}}\right] {\left({\ {\bmi  I}_{ 3 N_k} -
{\bmi S}{}_{\downarrow\uparrow }^{({s} , {i-1} ]}{\bmi S}{}_{\uparrow 
\downarrow
}^{( {i-1} , \infty )} }\right)}^{-1} {\bmi S}{}_{\downarrow
\downarrow }^{( {s} , {i-1} ]}
}\nonumber\\ 
&&\qquad \qquad \qquad \times{\left({ {\bmi  I}_{ 3 N_k} - 
{\bmi S}{}_{\downarrow \uparrow }^{[ f , {s} )} {\bmi S}{}_{\uparrow
\downarrow }^{( {s} , \infty )} }\right)}^{-1}
\left[{\bmi S}{}_{\downarrow \uparrow }^{[f, {s})}{\bmi \Sigma}_{\uparrow}+ {\bmi 
\Sigma}_{\downarrow} \right]
\end{eqnarray}
 Alternatively, for a source in the lower halfspace, the incident field in layer-$i$ is 
 given by 
\begin{eqnarray} \label{incup}
\lefteqn{
{\bmi  V}_{{inc}}^{ i}=\left[{\begin{array}{c}{\bmi  V}_{{inc},\downarrow}^{ 
i}\\ {\bmi  V}_{{inc},\uparrow}^{ i}\end{array}}\right]
=\left[{\begin{array}{c}
{\bmi S}{}_{\downarrow\uparrow }^{   [f,{i-1}]}{{\bmi E}}_{i}\\ 
{\bmi  I}_{3 N_k}\end{array}}\right] 
{\left({\bmi  I}_{ 3 N_k}-
{\bmi S}{}_{\uparrow\downarrow }^{ [i,{s})} {\bmi S}{}_{\downarrow\uparrow 
}^{ 
[f,i)}\right)}^{-1}\
 {\bmi S}{}_{\uparrow\uparrow }^{ [i,z_{s})}
}\nonumber\\
&&\qquad \qquad \qquad \times
{\left({{\bmi  I}_{3 N_k}-{\bmi S}{}_{\uparrow\downarrow }^{ ({s},\infty)}
{{\bmi S}}{}_{\downarrow\uparrow }^{ [f,{s})}}\right)}^{-1}
\left[ {\bmi \Sigma}_{\uparrow} +
{\bmi S}{}_{\uparrow\downarrow }^{({s}, \infty)} {\bmi \Sigma}_{\downarrow}
\right].
\end{eqnarray}






If we now model 
\[
{\bmi \Sigma} \equiv\Delta {\CB  S}{\CB  V}_{{int}}
\] 
in (\ref{LSmat}) as a {\em virtual \/} or {\em secondary\/} source 
field, then the the {\em scattered field\/} ${\mit \Delta}{\CB V}$ is given by
\begin{equation} \label{dSdef}
{\mit \Delta}{\CB V}=
({\bmi  I}-{\CB  S}_{{int}}{\CB  E})^{-1}{\CB  S}_{{int}} {\mit \Delta} 
{\CB  S}{\CB  V}_{{int}}=
{\CB Q}^{\partial}{\bmi \Sigma}.
\end{equation}
Partitioning the scattered field ${\mit \Delta}  { \CB V}$ and secondary source field
${\bmi \Sigma}$ in (\ref{dSdef}) according to
\begin{equation} \label{delVandSig}
 {\mit \Delta}  { \CB V}=\left[{\begin{array}{c}\left[{\begin{array}{c} 
{\mit \Delta} {  \bmi V}_{\downarrow }^{  1}\\ {\mit \Delta} {  
\bmi V}_{\uparrow }^{  1}\end{array}}\right]\\ \vdots\\ 
\left[{\begin{array}{c} {\mit \Delta} {  \bmi V}_{\downarrow }^{  
L}\\ {\mit \Delta} {\bmi  V}_{\uparrow }^{  
L}\end{array}}\right]\end{array}}\right] , \quad \mbox{and} \quad
 {\bmi \Sigma}=\left[{\begin{array}{c}\left[{\begin{array}{c}
{\bmi \Sigma}_{\uparrow }^{  1}\\
{\bmi \Sigma}_{\downarrow }^{  1}\end{array}}\right]\\
\vdots\\
\left[{\begin{array}{c} {\bmi \Sigma}_{\uparrow }^{  L}\\
{\bmi \Sigma}_{\downarrow }^{  L}\end{array}}\right]\end{array}}\right]
\end{equation}
then,  (\ref{dSdef}) may be written as  (Fig.\ref{Qfig})
 \begin{eqnarray} \label{superPdef}
\lefteqn{
\left[{\begin{array}{c}\left[{\begin{array}{c} 
{\mit \Delta} {  \bmi V}_{\downarrow }^{  1}\\ {\mit \Delta} {  
\bmi V}_{\uparrow }^{  1}\end{array}}\right]\\ \vdots\\ 
\left[{\begin{array}{c} {\mit \Delta} {  \bmi V}_{\downarrow }^{  
L}\\ {\mit \Delta} {\bmi  V}_{\uparrow }^{  
L}\end{array}}\right]\end{array}}\right] =
	\left({\begin{array}{ccc}
\left[{\begin{array}{cc}
({\bmi Q}^{\partial}_{1\leftarrow 1})_{11}&
({\bmi Q}^{\partial}_{1 \leftarrow 1})_{12}\\ 
({\bmi Q}^{\partial}_{1 \leftarrow  1})_{21}&
({\bmi Q}^{\partial}_{1 \leftarrow 1})_{22}\end{array}}\right]
& \cdots &
\left[{\begin{array}{cc}
({\bmi Q}^{\partial}_{1\leftarrow L})_{11}&
({\bmi Q}^{\partial}_{1 \leftarrow L})_{12}\\ 
({\bmi Q}^{\partial}_{1 \leftarrow  L})_{21}&
({\bmi Q}^{\partial}_{1 \leftarrow L})_{22}\end{array}}\right]\\ 
\vdots&\cdots &\vdots\\ 
\left[{\begin{array}{cc}
({\bmi Q}^{\partial}_{L\leftarrow 1})_{11}&
({\bmi Q}^{\partial}_{L \leftarrow 1})_{12}\\ 
({\bmi Q}^{\partial}_{L \leftarrow  1})_{21}&
({\bmi Q}^{\partial}_{L \leftarrow 1})_{22}\end{array}}\right]
 & \cdots &
\left[{\begin{array}{cc}
({\bmi Q}^{\partial}_{L\leftarrow L})_{11}&
({\bmi Q}^{\partial}_{L \leftarrow L})_{12}\\ 
({\bmi Q}^{\partial}_{L \leftarrow  L})_{21}&
({\bmi Q}^{\partial}_{L \leftarrow L})_{22}\end{array}}\right]
 \end{array}}\right)
 \left[{\begin{array}{c}\left[{\begin{array}{c}
{\bmi \Sigma}_{\uparrow }^{  1}\\
{\bmi \Sigma}_{\downarrow }^{  1}\end{array}}\right]\\
\vdots\\
\left[{\begin{array}{c} 
{\bmi \Sigma}_{\uparrow }^{  L}\\
{\bmi \Sigma}_{\downarrow }^{  
L}\end{array}}\right]\end{array}}\right]
}\nonumber\\
&&
\end{eqnarray}
where an explicit construction for ${\bmi Q}^{\partial}$ is presented in {\sc Box} {\bf 5.4}
%%%%%%%%%%%%%%%%%%%%%%%%%%%%%%%%%%%%%%%%%%%%%%%%%%%%%%%%%%%%%%%%%%%%%%%%%%%%%%%%%%%%%%%%%%%
\begin{figure}
\HideDisplacementBoxes
%\ShowDisplacementBoxes
\centerline{\BoxedEPSF{Q.EPSF scaled 1000}}
\caption{
The {\em plane wave propagator\/} 
$ {\mbox{\mib Q}}^{\partial}_{i \leftarrow j}$ 
maps the {\em virtual source field\/} ${\bf \Sigma}^{ j}\equiv (
{\bf \Sigma}_{\uparrow}^{ j}, {\bf \Sigma}_{\downarrow}^{ j})^T$ from layer-$j$ to the 
{\em scattered field\/} ${\mit \Delta}\mbox{\mib V}^{ i}\equiv (
{\mit \Delta}\mbox{\mib V}_{\downarrow}^{ i}, {\mit \Delta}\mbox{\mib V}_{\uparrow}^{ i})^T$  in layer-$i$.
}
\protect\label{Qfig}
\end{figure}
%%%%%%%%%%%%%%%%%%%%%%%%%%%%%%%%%%%%%%%%%%%%%%%%%%%%%%%%%%%%%%%%%%%%%%%%%%%%%%%%%%%%%%%%%%%
 
\begin{samepage} 
\begin{figure}
\[
\!\!\!\!\!\!\!\!\!\!\!\!\!\!\!\!\!\!\!\!\!\!
\fbox{
\begin{minipage} {160.0mm}
{\sc \bf  Box} $\bf 5.4$ 
{\sc An explicit representation of the ${\bmi Q}^{\partial}$ propagator\/} \newline
\vspace{4mm}
\begin{eqnarray*} 
 \lefteqn{ 
 {\bmi Q}^{\partial}_{m\leftarrow  l}=
 \left[({\bmi  I}-{\CB  S}_{{int}}{\CB  E})^{-1}{\CB  
 S}_{{int}}\right]_{m,l}
 }\\
&&\!\!\!\!\!\!\!\!\!\!\!\!\!\!\!\!\!\!
\left\{
\begin{array}{ll}
\left[{\begin{array}{c}{{\bmi I}}\\ 
{\bmi S}{}_{\uparrow\downarrow }^{(m-, 
L)}{{\bmi E}}_{m}\end{array}}\right] {\left({\ {\bmi  I}_{ 3 N} -
{\bmi S}{}_{\downarrow\uparrow }^{( l- , {m-1}+ )}{\bmi S}{}_{\uparrow 
\downarrow
}^{( {m-1}+ , L )} }\right)}^{-1} {\bmi S}{}_{\downarrow
\downarrow }^{( l- , {m-1}+ )}
&\\ 
\qquad \qquad \qquad \qquad \qquad \qquad \qquad \times{\left({ {\bmi  I}_{ 3 N} - 
{\bmi S}{}_{\downarrow \uparrow }^{( f , l- )} {\bmi S}{}_{\uparrow
\downarrow }^{( l- , L )} }\right)}^{-1}
\left[{\begin{array}{cc}{{\bmi E}}_{l}{\bmi S}{}_{\downarrow \uparrow }^{(f, 
l-1+)}&{{\bmi I}}\end{array}}\right]&\!\!\!:m>l\\
&\\
\!\!\!\!\!
\left[\begin{array}{cc}
{\left({ {\bmi  I} - {\bmi S}{}_{\downarrow \uparrow }^{(f,
{m-1})}{\bmi S}{}_{\uparrow\downarrow }^{({m-1}, L)} }\right)}^{-1}
{\bmi S}{}_{\downarrow \uparrow }^{(f, {m-1})}&{\bmi S}{}_{\downarrow\uparrow }^{(f, {m-1})} 
{\left({ {\bmi  I}- {\bmi S}{}_{\uparrow \downarrow}^{({m-1},
L)}{\bmi S}{}_{\downarrow\uparrow }^{(f, {m-1})} }\right)}^{-1}
{{\bmi E}}_{ m} {\bmi S}_{\uparrow \downarrow }^{
(m, L)}\\\!
{\bmi S}{}_{\uparrow \downarrow }^{(m, L)} 
{{\bmi E}}_{ 
m} 
{\left({   {\bmi  I} -{\bmi S}{}_{\downarrow
\uparrow }^{(f, {m-1})}{\bmi S}{}_{\uparrow
\downarrow }^{({m-1}, L)} }\right)}^{  -1} {\bmi S}{}_{\downarrow 
\uparrow 
}^{(f, {m-1})}&{\bmi S}{}_{\uparrow\downarrow }^{(m, L)} {\left({ 
{\bmi I} - {\bmi S}{}_{\downarrow \uparrow }^{(f,
m)} {\bmi S}{}_{\uparrow\downarrow }^{(m, L)} }\right)}^{-1}
\end{array}\!\right]&\!\!\!:m=l\\
&\\
\left[{\begin{array}{c}
{\bmi S}{}_{\downarrow\uparrow }^{   (f,{m-1}+)}{{\bmi E}}_{m}\\ 
{{\bmi I}}\end{array}}\right] 
{\left({\bmi  I}_{ 3 N}-
{\bmi S}{}_{\uparrow\downarrow }^{ (m-,L+)} {\bmi S}{}_{\downarrow\uparrow 
}^{(f,m-)}\right)}^{-1}\
 {\bmi S}{}_{\uparrow\uparrow }^{ (m-,{l-1}+)}\\
\qquad \qquad \qquad \qquad \qquad \qquad \qquad \times
{\left({{\bmi  I}_{
3 N}-{\bmi S}{}_{\uparrow\downarrow }^{ ({l-1}+,L+)}
{{\bmi S}}{}_{\downarrow\uparrow }^{ (m-,{l-1}+)}}\right)}^{-1}
\left[{\begin{array}{cc}{{\bmi I}}&{{\bmi E}}_{l}{\bmi 
S}{}_{\uparrow\downarrow }^{(l-, 
L+)}\end{array}}\right]&\!\!\!:m<l
\end{array}\right.
\end{eqnarray*}

\vspace{11pt}

\end{minipage}
}
\]
\end{figure}
\end{samepage}
 
The {\em matrix Lippmann-Schwinger equation\/} (\ref{LSmat}) may now be written in the 
slightly more abstract form
\begin{equation} \label{LSmat1}
{\CB  V}_{{int}}={\CB  V}_{{inc}} + {\mit \Delta}  { \CB V}
\end{equation}
so that the {\em total field\/} ${\CB  V}_{{int}}$, {\em phased relative to the reference medium 
interfaces\/},  is the sum of the {\em incident 
field\/} ${\CB  V}_{{inc}}$ and the {\em scattered field\/} ${\mit \Delta}  
{ \CB V}$ from the superimposed heterogeneity.



Due to the enormous matrix dimension, we must be very careful to optimize the efficiency of the matrix 
operation (\ref{dSdef}). We now propose a {\em Riccati matrix factorization 
scheme\/} based on {\em reflection/transmission matrix recursions\/} that 
are quite similar to those of {\em Kennett\/}~(1983), Chapter 
\ref{exten} (indeed, we 
make use of the Kennett's recursions in our derivation), but are 
more convenient for the {\em multi-source, multi-receiver\/} application, 
currently under investigation.   
To begin the derivation of the Riccati factorization, we observe 
that the array (\ref{mollref}) is {\em block tri-diagonal\/}, i.e.
\begin{equation} \label{Amat}
{\CB  A} \equiv ({\bmi  I}-{\CB  S}_{{int}}{\CB  E}) = 
\mbox{tridiag}\!\! \left\{{\begin{array}{llll}{\bmi  A}_{0,0}&{\bmi 
A}_{0,1}&&i=0\\
{\bmi  A}_{i,i-1}&{\bmi  A}_{i,i}&{\bmi  A}_{i,i+1}&i=1,\cdots ,L-1\\
&{\bmi  A}_{L,L-1}&{\bmi  A}_{L,L}&i=L\end{array}}\right. ;
\end{equation}
where the matrix partitions of (\ref{Amat}) are defined:
\begin{equation} 
{\bmi  A}_{ 0, 0}={\bmi  A}_{ L, L}= {\bmi  I}_{3 N_k} , 
\ {\bmi  A}_{ 0, 1}=
\left(\begin{array}{c} -{\bmi S}^{[f,0]}_{\downarrow\uparrow 
} {\bmi  E}_{1}\end{array}\right),
\
{\bmi  A}_{  L,L-1}=
\left(\begin{array}{c} 
-{\bmi S}^{[L,\infty)}_{\uparrow\downarrow} {\bmi  E}_{L}\end{array}\right)
\end{equation}
and the remaining elements of $\CB  A$ are $6 N_k\times 6 N_k$ arrays
\begin{equation}
\!\!\! 
{\bmi  A}_{ i, i}= {\bmi  I}_{6 N_k} ,
{\bmi  A}_{i, i-1} =
\left({\begin{array}{cc} {\bf 0}_{3 N_k}&-{\bmi S}^{i}_{\uparrow 
\downarrow} {\bmi  E}_{i}\\ 
{\bf 0}_{3 N_k}&-{\bmi S}^{i}_{\downarrow\downarrow } {\bmi  E}_{i}
\end{array}}\right) , 
\ 
{\bmi  A}_{i, i+1}= \left({\begin{array}{cc} 
-{\bmi S}^{i}_{\uparrow \uparrow
} {\bmi  E}_{i+1}&{\bf 0}_{3 N_k}\\ -{\bmi S}^{i}_{\downarrow\uparrow 
} {\bmi  E}_{i+1}&{\bf 0}_{3 N_k}\end{array}}\right) 
\end{equation}
except for the $6 N_k\times 3 N_k$ arrays
\begin{equation}	\label{A10def}
	{\bmi  A}_{ 1, 0}=\left({\begin{array}{c} -{\bmi S}^{1}_{\uparrow 
\downarrow } {\bmi  E}_{1}\\ 
	-{\bmi S}^{1}_{\downarrow\downarrow } {\bmi  E}_{1}\end{array}}\right), \
{\bmi  A}_{L-1, L}= \left({\begin{array}{c} -{\bmi S}^{L-1}_{\uparrow 
\uparrow} {\bmi  E}_{L}\\ -{\bmi S}^{L-1}_{\downarrow\uparrow } {\bmi 
E}_{L}
\end{array}}\right).
\end{equation}

Although the matrix dimension of the reference 
medium reflection/transmission arrays such as 
${\bmi S}{}_{\uparrow\downarrow }^{i}$ in (\ref{Amat})  is very large in general 
($\approx 10^{4}$ in the applications considered in this work), 
{\em due to the lateral invariance of the reference medium, they are quasi-diagonal\/} {\em 
cf.\/}(\ref{Ssparsity}).
We can exploit the block tridiagonal structure of ${\CB  A}$ and the {\em 
sparsity\/} (\ref{Ssparsity}) of the components, to calculate 
the action of ${\CB  A}^{-1}$,
on a field vector, via the factorization 
\begin{eqnarray}
\lefteqn{{\CB  A} ={\CB  U}{\CB  L}=}\nonumber\\
&&\!\!\!\!\!\!
\left[{\begin{array}{cc}\begin{array}{ccc}{\bmi  D}^{-1}_{0}&{\bmi 
A}_{01}&\\ &{\bmi  D}^{-1}_{1}&{\bmi  A}_{12}\\&&\ddots\end{array}&\\
&\begin{array}{cc}{\bmi  D}^{-1}_{L-1}&{\bmi  A}_{ L-1,L}\\ &{\bmi  D}^{-1}_{L}
\end{array}\end{array}}\right]
\left[{\begin{array}{cc}
\begin{array}{ccc}{\bmi  I}_{3 N_k}&&\\{\bmi 
D}_{1}{\bmi  A}_{10} &{\bmi  I}_{6 N_k}&\\&&\ddots\end{array}&\\ 
&\begin{array}{cc}
{\bmi  I}_{6 N_k}&\\ {\bmi  D}_{L}{\bmi  A}_{L,L-1}&{\bmi 
I}_{3 N_k}\end{array}\end{array}}\right]
\end{eqnarray}
where the diagonal blocks ${\bmi  D}_{i}$ are constructed according to
\begin{eqnarray} \label{Dblock}
{\bmi  D}_{i}&=&\left({\bmi  A}_{i,i}- {\bmi  A}_{i,i+1}{\bmi  D}_{i+1}{\bmi 
A}_{i+1,i}
\right)^{-1}\\
&=&\left\{{\begin{array}{cl}{\left({{\bmi  I}_{3 N_k}-{\bmi S}_{ \downarrow 
\uparrow 
}^{[f,0]}{\bmi S}_{ \uparrow \downarrow }^{(0,\infty)}}\right)}^{-1}& \qquad:\qquad
i=0\\ 
\left({\begin{array}{cc}{\bmi  I}_{3 N_k}&{{\bmi S}_{ \uparrow \uparrow 
}^{i}{\bmi S}_{ 
\uparrow \downarrow }^{(i,\infty)}\left({{\bmi  I}_{3 N_k}-{\bmi S}_{ 
\downarrow \uparrow }^{i}{\bmi S}_{ \uparrow \downarrow 
}^{(i,\infty)}}\right)}^{-1}\\ 
{\bf 0}_{3 N_k}&{\left({{\bmi  I}_{3 N_k}-{\bmi S}_{ \downarrow \uparrow 
}^{i}{\bmi S}_{ \uparrow \downarrow }^{(i,\infty)}}\right)}^{-1}\end{array}}\right)& \qquad:\qquad
i=1,\cdots ,L-1\\ {\bmi  I}_{3 N_k}&\qquad:\qquad i=L\end{array}}\right.\nonumber
\end{eqnarray}
where ${\bmi S}_{ \uparrow \downarrow}^{(i,\infty)}$ in (\ref{Dblock}) 
is the {\em net downward reflection operator\/} for the region of the layer 
stack below interface-$i$ (i.e. $z_{i} < z < \infty$), and is defined recursively in terms of 
${\bmi S}_{ \uparrow \downarrow}^{(i+1, \infty)}$ as 
\begin{equation} \label{Rrecursion}
{\bmi S}_{\uparrow \downarrow }^{(i,\infty)}={\bmi E}_{i+1}
\left[{{\bmi S}_{ \uparrow \downarrow }^{i+1}+{\bmi S}_{ \uparrow 
\uparrow }^{i+1}{\bmi S}_{ \uparrow \downarrow }^{(i+1, \infty   
)}\left({{\bmi I }-{\bmi S}_{ \downarrow \uparrow }^{i+1}{\bmi S}_{ \uparrow \downarrow 
}^{(i+1, \infty   )}}\right){\bmi S}_{\downarrow 
\downarrow}^{i+1}}\right]{\bmi E}_{i+1}.
\end{equation} 
Eqn. (\ref{Rrecursion}) is essentially the  recursion for the net downward 
reflection operator [{\em Kennett\/}~(1983) Eqn. (6.3)]. 

%%%%%%%%%%%%%%%%%%%%%%%%%%%%%%%%%%%%%%%%%%%%%%%%%%%%%%%%%%%%%%%%%%%%%%%%%%%%%%%%%%%%%%%%%%%
%%%%%%%%%%%%%%%%%%%%%%%%%%%%%%%%%%%%%%%%%%%%%%%%%%%%%%%%%%%%%%%%%%%%%%%%%%%%%%%%%%%%%%%%%%%
%%%%%%%%%%%%%%%%%%%%%%%%%%%%%%%%%%%%%%%%%%%%%%%%%%%%%%%%%%%%%%%%%%%%%%%%%%%%%%%%%%%%%%%%%%%
\begin{figure}
%\vspace*{10em}
\fbox{
\begin{minipage} {140.0mm}
{\bf  Algorithm}~{\bf 5.1} {\sc A Riccati sweep implementation of\/} 
${\mit \Delta}{\CB V}=({\bmi  I}-{\CB  S}_{{int}}{\CB  E})^{-1}{\CB  S}_{{int}} {\bmi \Sigma}
={\CB Q}^{\partial}{\bmi \Sigma}$\newline
\vspace{4mm}

\begin{itemize}
\item given the {\em virtual source field\/} $\bmi \Sigma$ compute the 
{\em scattered field\/} ${\mit \Delta}{\CB V}
={\CB Q}^{\partial}{\bmi \Sigma}$

\begin{itemize}
\item  (*{\em bottom of layer stack\/ }*) 
\end{itemize}
\item ${\bmi R}_{ \uparrow \downarrow }^{L}\equiv 
{\bmi S}_{ \uparrow\downarrow }^{[L,\infty)}$
\item  	$
{\bmi t}_{ \uparrow }^{L}=
{\bmi R}_{ \uparrow \downarrow }^{L}
{\bmi \Sigma}^{L}_{\downarrow}$

\begin{itemize}
\item  (*{\em upsweep\/ }*) 
\end{itemize}
\item {\bf  For} $i=L-1,\cdots,1$
\begin{itemize}

\item ${\bmi R}_{ \uparrow \downarrow }^{i}=
{\bmi S}_{ \uparrow \downarrow }^{i}+{\bmi S}_{ \uparrow \uparrow 
}^{i}
{\bmi  E}_{i+1} {\bmi R}_{ \uparrow \downarrow }^{i+1} {\bmi  E}_{i+1} {\left({{\bmi  I}-{\bmi S}_{ \downarrow 
\uparrow }^{i}
{\bmi  E}_{i+1} {\bmi R}_{ \uparrow \downarrow }^{i+1} {\bmi  E}_{i+1} }\right)}^{-1}{\bmi S}_{ \downarrow 
\downarrow }^{i}
 $

\item ${\bmi R}_{ \uparrow \uparrow }^{i}= {\bmi S}_{ \uparrow \uparrow 
}^{i}
{\left({{\bmi  I}-{\bmi  E}_{i+1} {\bmi R}_{ \uparrow \downarrow }^{i+1} {\bmi 
E}_{i+1} {\bmi S}_{ \downarrow 
\uparrow }^{i}}\right)}^{-1}$

\item ${\bmi R}_{ \downarrow \downarrow }^{i}={\left({{\bmi  I}-{\bmi S}_{ 
\downarrow \uparrow }^{i} {\bmi  E}_{i+1} {\bmi R}_{ \uparrow \downarrow 
}^{i+1} {\bmi  E}_{i+1} }
\right)}^{-1}{\bmi S}_{ \downarrow \downarrow }^{i} $

\item ${\bmi R}_{ \downarrow \uparrow }^{i} = 
{\left({{\bmi  I}-{\bmi S}_{ \downarrow \uparrow }^{i}
{\bmi  E}_{i+1} {\bmi R}_{ \uparrow \downarrow }^{i+1} {\bmi  E}_{i+1} }\right)}^{-1}
{\bmi S}_{ \downarrow \uparrow }^{i}$



\item  $\left[{\begin{array}{c}
{\bmi t}_{ \uparrow }^{i}\\ {\bmi t}_{ \downarrow }^{i+1}\end{array}}\right]=
\left[{\begin{array}{cc}
{\bmi R}_{ \uparrow \downarrow }^{i}&{\bmi R}_{ \uparrow \uparrow }^{i}\\ 
{\bmi R}_{ \downarrow \downarrow }^{i}&{\bmi R}_{ \downarrow \uparrow 
}^{i}\end{array}}\right]
\left[{\begin{array}{c}{\bmi \Sigma}^{i}_{\downarrow}\\ 
{\bmi \Sigma}^{i+1}_{\uparrow} + {\bmi  E}_{i+1}{  \bmi t}_{ \uparrow }^{i+1}
\end{array}}\right]$
\end{itemize}
\item {\bf  End}

\begin{itemize}
\item  (*{\em top of layer stack\/ }*) 
\end{itemize}
\item ${\bmi R}_{ \downarrow \uparrow }^{0} = 
{\left({{\bmi  I}-{\bmi S}_{ \downarrow \uparrow }^{[f,{0}]}
{\bmi  E}_{1} {\bmi R}_{ \uparrow \downarrow }^{1} {\bmi  E}_{1} }\right)}^{-1}
{\bmi S}_{ \downarrow \uparrow }^{[f,{0}]}$


\item  ${\bmi t}_{ \downarrow }^{1}={\bmi R}_{ \downarrow \uparrow }^{0}
({\bmi \Sigma}^{1}_{\uparrow}+{\bmi  E}_{i}{  \bmi t}_{ \uparrow }^{1})$

\item  ${\mit \Delta}{{\bmi V}}_{ \downarrow }^{1}={\bmi t}_{ \downarrow }^{1}$

\begin{itemize}
\item  (*{\em downsweep\/ }*) 
\end{itemize}
\item {\bf  For} $i=1,\cdots,L-1$
\begin{itemize}
\item  $\left[{\begin{array}{c}
{\mit \Delta}{{\bmi V}}_{ \uparrow }^{i}\\ {\mit \Delta}{{\bmi V}}_{ 
\downarrow }^{i+1}\end{array}}\right]=\left[{\begin{array}{c}{  \bmi R}_{ 
\uparrow \downarrow }^{i}\\ {  \bmi R}_{ \downarrow \downarrow 
}^{i}\end{array}}\right] {\bmi  E}_{i} {  {\bmi V}}_{ \downarrow 
}^{i}+\left[{\begin{array}{c}{  \bmi t}_{ \uparrow }^{i}\\ {  
\bmi t}_{ \downarrow }^{i+1}\end{array}}\right] $
\end{itemize}
\item {\bf  End}


\begin{itemize}
\item  (*{\em bottom of layer stack\/ }*) 
\end{itemize}
\item  	$
{\mit \Delta}{{\bmi V}}_{ \uparrow }^{L}=
{\bmi R}_{ \uparrow \downarrow }^{L}
{\bmi  E}_{L} {{\bmi V}}_{ \downarrow }^{L} +{\bmi t}_{ \uparrow }^{L}$


\end{itemize}

\vspace{6 mm}
\end{minipage}
}
\end{figure}
%%%%%%%%%%%%%%%%%%%%%%%%%%%%%%%%%%%%%%%%%%%%%%%%%%%%%%%%%%%%%%%%%%%%%%%%%%%%%%%%%%%%%%%%%%%
%%%%%%%%%%%%%%%%%%%%%%%%%%%%%%%%%%%%%%%%%%%%%%%%%%%%%%%%%%%%%%%%%%%%%%%%%%%%%%%%%%%%%%%%%%%
%%%%%%%%%%%%%%%%%%%%%%%%%%%%%%%%%%%%%%%%%%%%%%%%%%%%%%%%%%%%%%%%%%%%%%%%%%%%%%%%%%%%%%%%%%%
%%%%%%%%%%%%%%%%%%%%%%%%%%%%%%%%%%%%%%%%%%%%%%%%%%%%%%%%%%%%%%%%%%%%%%%%%%%%%%%%%%%%%%%%%%%


We now construct the scattered field ${\mit \Delta}{\CB V}$ in (\ref{dSdef}) by 
solving the triangular systems
\begin{equation} \label{ULsweeps}
\begin{array}{l}
{\CB U}{\bmi t}={\CB S}_{{int}}{\bmi \Sigma}\\
{\CB L}{\mit \Delta}{\CB V}={\bmi t}
\end{array}
\end{equation}
where:\newline
the right hand side vector ${\CB S}_{{int}}{\bmi \Sigma}$ in 
(\ref{ULsweeps}) is defined
\begin{eqnarray}
{\CB S}_{{int}}{\bmi \Sigma}&=&
\left[{\begin{array}{ccccc}
{\bmi S}_{\downarrow\uparrow}^{[f,0]}&&&&\\ 
&\begin{array}{cc}{\bmi S}_{\uparrow
\downarrow }^{1}&{\bmi S}_{\uparrow \uparrow }^{1}\\ 
{\bmi S}_{\downarrow
\downarrow }^{1}&{\bmi S}_{\downarrow \uparrow 
}^{1}\end{array}&&&\\ && \ddots
&&\\  &&&\begin{array}{cc}{\bmi S}_{\uparrow \downarrow 
}^{L-1}&{\bmi S}_{\uparrow \uparrow }^{L-1}\\ 
{\bmi S}_{\downarrow \downarrow 
}^{L-1}&{\bmi S}_{\downarrow \uparrow 
}^{L-1}\end{array}&\\ 
&&&&{\bmi S}_{\uparrow\downarrow}^{[L,\infty)}\end{array}}\right]
\left[{\begin{array}{c}{\bmi \Sigma}_{ \uparrow }^{1}\\
\left[{\begin{array}{c}{\bmi \Sigma}_{ \downarrow }^{1}\\
{\bmi \Sigma}_{ \uparrow }^{2}\end{array}}\right]\\
\vdots\\
\left[{\begin{array}{c}{\bmi \Sigma}_{ \downarrow }^{L-1}\\
{\bmi \Sigma}_{ \uparrow }^{L}\end{array}}\right]\\
{\bmi \Sigma}_{ \downarrow }^{L}
\end{array}}\right]\nonumber\\
&=&
\left[{\begin{array}{c}
{\bmi S}_{ \downarrow \uparrow }^{[f,0]}{ 
\bmi \Sigma}_{ \uparrow }^{1}\\ \left[{\begin{array}{c}{\bmi S}_{ \uparrow 
\downarrow }^{1}{ \bmi \Sigma}_{ \downarrow }^{1}+{\bmi S}_{ \uparrow \uparrow 
}^{1}{ \bmi \Sigma}_{ \uparrow }^{2}\\ {\bmi S}_{ \downarrow \downarrow 
}^{1}{ \bmi \Sigma}_{ \downarrow }^{1}+{\bmi S}_{ \downarrow \uparrow 
}^{1}{ \bmi \Sigma}_{ \uparrow }^{2}\end{array}}\right]\\ \vdots\\ 
\left[{\begin{array}{c}{\bmi S}_{ \uparrow \downarrow }^{L-1}{ \bmi \Sigma}_{ 
\downarrow }^{L-1}+{\bmi S}_{ \uparrow \uparrow }^{L-1}{ \bmi \Sigma}_{ 
\uparrow }^{L}\\ {\bmi S}_{ \downarrow \downarrow }^{L-1}{ \bmi \Sigma}_{ 
\downarrow }^{L}+{\bmi S}_{ \downarrow \uparrow }^{L-1}{ \bmi \Sigma}_{ 
\uparrow }^{L}\end{array}}\right]\\ {\bmi S}_{ \uparrow \downarrow }^{[L, 
\infty   )}{ \bmi \Sigma}_{ \downarrow }^{L}\end{array}}\right] \nonumber;
\end{eqnarray}
we solve the {\em lower triangular system\/}
\begin{samepage}
\begin{eqnarray} \label{Upsweep}
\left[{\begin{array}{cc}\begin{array}{ccc}{\bmi  D}^{-1}_{0}&{\bmi 
A}_{01}&\\ &{\bmi  D}^{-1}_{1}&{\bmi  A}_{12}\\&&\ddots\end{array}&\\
&\begin{array}{cc}{\bmi  D}^{-1}_{L-1}&{\bmi  A}_{ L-1,L}\\ &{\bmi  D}^{-1}_{L}
\end{array}\end{array}}\right]
\left[{\begin{array}{c}
{\bmi t}_{ \downarrow }^{1}\\
\left[{\begin{array}{c}
{\bmi t}_{ \uparrow }^{1}\\
{\bmi t}_{ \downarrow }^{2}\end{array}}\right]\\
\vdots\\
\left[{\begin{array}{c}
{\bmi t}_{ \uparrow }^{L-1}\\
{\bmi t}_{ \downarrow }^{L}\end{array}}\right]\\
{\bmi t}_{ \uparrow }^{L}\end{array}}\right]
&=&
\left[{\begin{array}{c}
{\bmi S}_{ \downarrow \uparrow }^{[f,0]}{ 
\bmi \Sigma}_{ \uparrow }^{1}\\ \left[{\begin{array}{c}{\bmi S}_{ \uparrow 
\downarrow }^{1}{ \bmi \Sigma}_{ \downarrow }^{1}+{\bmi S}_{ \uparrow \uparrow 
}^{1}{ \bmi \Sigma}_{ \uparrow }^{2}\\ {\bmi S}_{ \downarrow \downarrow 
}^{1}{ \bmi \Sigma}_{ \downarrow }^{1}+{\bmi S}_{ \downarrow \uparrow 
}^{1}{ \bmi \Sigma}_{ \uparrow }^{2}\end{array}}\right]\\ \vdots\\ 
\left[{\begin{array}{c}{\bmi S}_{ \uparrow \downarrow }^{L-1}{ \bmi \Sigma}_{ 
\downarrow }^{L-1}+{\bmi S}_{ \uparrow \uparrow }^{L-1}{ \bmi \Sigma}_{ 
\uparrow }^{L}\\ {\bmi S}_{ \downarrow \downarrow }^{L-1}{ \bmi \Sigma}_{ 
\downarrow }^{L}+{\bmi S}_{ \downarrow \uparrow }^{L-1}{ \bmi \Sigma}_{ 
\uparrow }^{L}\end{array}}\right]\\ {\bmi S}_{ \uparrow \downarrow }^{[L, 
\infty   )}{ \bmi \Sigma}_{ \downarrow }^{L}\end{array}}\right]\nonumber\\
&&
\end{eqnarray}
\end{samepage}
of (\ref{ULsweeps}), using {\em back substitution\/}, and we solve the {\em upper triangular\/}
system 
\begin{equation} \label{Downsweep}
\left[{\begin{array}{cc}\begin{array}{ccc}{\bmi  I}_{3 N_k}&&\\{\bmi 
D}_{1}{\bmi  A}_{10} &{\bmi  I}_{6 N_k}&\\&&\ddots\end{array}&\\ 
&\begin{array}{cc}
{\bmi  I}_{6 N_k}&\\ {\bmi  D}_{L}{\bmi  A}_{L,L-1}&{\bmi 
I}_{3 N_k}\end{array}\end{array}}\right]
\left[{\begin{array}{c} {\mit \Delta} {  
\bmi V}_{\downarrow }^{  1}\\ \left[{\begin{array}{c}{\mit \Delta} {  
\bmi V}_{\uparrow }^{  1}\\ {\mit \Delta} {  \bmi V}_{\downarrow }^{  
2}\end{array}}\right]\\   \vdots\\ \left[{\begin{array}{c} 
{\mit \Delta} {  \bmi V}_{\uparrow }^{  L-1}\\ {\mit \Delta} {  
\bmi V}_{\downarrow }^{  L}\end{array}}\right]\\  {\mit \Delta} {  
\bmi V}_{\uparrow }^{  L}\end{array}}\right]
=
\left[{\begin{array}{c}
{\bmi t}_{ \downarrow }^{1}\\
\left[{\begin{array}{c}
{\bmi t}_{ \uparrow }^{1}\\
{\bmi t}_{ \downarrow }^{2}\end{array}}\right]\\
\vdots\\
\left[{\begin{array}{c}
{\bmi t}_{ \uparrow }^{L-1}\\
{\bmi t}_{ \downarrow }^{L}\end{array}}\right]\\
{\bmi t}_{ \uparrow }^{L}\end{array}}\right]
\end{equation}
of (\ref{ULsweeps}), using {\em forward substitution\/}.
Also, we have introduced the auxiliary vector
\begin{equation} \label{tdef}
{\bmi t}=\left[{\begin{array}{ccccc}
{\bmi t}_{ \downarrow }^{1},&
\left[{\begin{array}{cc}
{\bmi t}_{ \uparrow }^{1},&
{\bmi t}_{ \downarrow }^{2}\end{array}}\right],&
\cdots,&
\left[{\begin{array}{cc}
{\bmi t}_{ \uparrow }^{L-1},
{\bmi t}_{ \downarrow }^{L}\end{array}}\right],&
{\bmi t}_{ \uparrow }^{L}\end{array}}\right]^{T}
\end{equation}
to save the result of the upsweep (\ref{Upsweep}). 

The sweeps (\ref{Upsweep}) and (\ref{Downsweep}) are summarized in  Alg.5.1.
This algorithm  involves the manipulation of a large number of reference 
medium reflection and transmission arrays. These calculations may be performed in an efficient manner, 
by direct manipulation of the component diagonal partitions; a technique 
that proved highly {\em vectorizable\/}, when implemented on a {\em Fujitsu\/} 
VP$2200$ supercomputer, at the Australian National University.
 










%%%%%%%%%%%%%%%%%%%%%%%%%%%%%%%%%%%%%%%%%%%%%%%%%%%%%%%%%%%%%%%%%%%%%%%%%%%%%%%%%%%%%
\subsection[forward scattering based splitting]
{A  forward scattering based splitting  and the two way 
marching algorithm}\label{fscat}

We now describe an alternative splitting of 
\[
({\bmi  I}-{\CB  S}_{{int}}{\CB  S}_{{layer}}){\CB  V}_{{int}}={\bmi 
b}_{{int}},
\]
from (\ref{schur}), that generalizes the {\em complex screen method\/} proposed by {\em 
Wu\/}~(1994). This splitting  is best suited to {\em smooth\/} reference 
media.

Modeling wave propagation in heterogeneous media is expensive. Indeed, 
although potentially complete,  discrete grid modeling schemes for the elastic wavefield such as 
finite difference-elements and pseudo-spectral methods are now well developed, the 
sheer volume of computation needed to solve realistic seismic wave 
problems, precludes them from routine use.
Born scattering has been proposed by several authors [e.g. {\em 
Kennett\/}~(1986)] as a reasonably cheap wave equation based alternative 
to the discrete grid methods. In the Born scattering philosophy,
the actual medium is cleaved into the sum of a reference (usually homogeneous or 
plane stratified), and a residual model. We then assume that the 
wavefield propagates through the reference medium, and the additional 
material variations of the actual medium that were neglected from the 
reference model serve as a continuum of independent {\em point scatterers\/}, 
modifying the phase and amplitude of the reference wave field via  {\em single 
scattering interactions}. 
 The consequence of this modeling strategy is that
``{\em the phases or travel times of both direct and scattered signals are 
habitually mistimed\/}" (Coates \& Charrette, 1993). The body of this 
thesis has concerned itself with re-introducing the neglected {\em 
multiple scattering\/} interactions, to correct for  this {\em mistiming}.

An alternative correction to Born scattering,  is the 
{\em elastic wave complex screen method\/} (Wu, 1994), where
the least time (or direct) wavefield propagates between successive forward scattering 
interactions at the correct wave speed. The mistiming, was originally 
corrected by applying ray theory (Fisk, Charrette, \& McCartor, 1992),
 which is of course a least time theory 
(being derived from the {\em least action\/} principle of Fermat). 
However, it has been observed by {\em Wu\/}~(1994), that 
the least time or direct arrivals may be modeled using wave theory also, 
by adding {\em the cumulative forward scattered field to the reference 
field, to get the wave theory direct field}. We now present a brief 
critique of the {\em elastic wave complex screen method\/}, because the 
techniques developed in this thesis may be applied to extend it into a 
{\em partially interacting wave theory}, i.e. we re-introduce the 
neglected backscatter. However, in contrast to the multiple scattering 
formulation of earlier chapters, here we employ the {\em thin slab Born 
approximation\/} of {\em Wu\/}~(1994), so that all  points at a fixed depth are {\em 
independent point scatterers, and do not reverberate amongst themselves} 
(i.e. {\em side scattering\/} is neglected). 


The {\em phase screen\/} method is based on two approximations:
\begin{itemize}
\item it exploits the assumed predominance of {\em forward scattering\/} over 
{\em back scattering \/} at high frequencies, so that to a first 
approximation, the wave field may 
be marched through the medium, by successive forward scattering 
interactions. The neglect of {\em back scattering\/}, avoids the need to 
invert the large reverberation arrays, such as 
\[
{\left[{\bmi I}-{\bmi S}_{\downarrow \uparrow }^{(A,B)}{\bmi S}_{\uparrow 
\downarrow }^{(B,C)}\right]}^{-1}
\] 
where ${\bmi S}_{\downarrow \uparrow }^{(A,B)}$ is the {\em net upward 
incident reflection array\/} for the depth interval $(A,B)$ above $z_B$ and ${\bmi S}_{\uparrow 
\downarrow }^{(B,C)}$ is the {\em net downward incident
reflection array\/} for the depth interval $(B,C)$ below $z_B$. Various combinations 
of these reverberation operators occur in the more complete theory based
on {\em coupled plane wave reflection-transmission operator recursions\/} 
(Kennett, 1984a; Kennett, 1984b).
\item the forward scattering interactions are modelled by passing the 
wavefield through a series of screens. The key property of the phase screens 
is that their effect on the wavefield is local in space (i.e. they have 
diagonal matrix representations in the spatial domain). As a consequence, 
the forward scattering interactions are  efficient to effect.
\end{itemize} 

In the phase screen method of {\em Fisk, Charrette, \& McCartor\/}~(1992),
the wave field  propagation between successive screens is based on ray theory.
The requirement of an  accurate model for P-SV wave coupling, has prompted 
{\em Wu\/}~(1994) to reformulate the ray theory based method of {\em Fisk 
et al.\/}~(1992), by applying a {\em  Born forward scattering method\/},
where the {\em cumulative forward scattered field is marched forward}. 
In this way, the incident field has the correct {\em least time phase\/} 
at each screen, 
correcting for the {\em habitual mis-timing\/} of the traditional Born method.
An alternate interpretation of the phase screen method  is: {\em if solving the multiple scattering equations using the Born approximation, 
is equivalent to a single iteration of the Jacobi method, then 
solving the same equations with the method proposed by\/} {\em 
Wu\/}~(1994) {\em
is equivalent to a single iteration of the  Gauss-Seidel method\/}. 
In the second stage of {\em Wu\/}~(1994),  
a {\em narrow angle\/} or {\em parabolic\/} approximation  to the 
forward scattering operators at each screen is proposed, so that the 
forward scattering operator may be approximated by a diagonal matrix in the spatial domain. 

There are two potential problems with {\em Wu\/}~(1994):
\begin{itemize}
\item  the {\em elastic wave 
complex screen method\/} is based on the premise that the net energy 
scattered forward, is larger than the net energy scattered backward for a 
medium extending over several wavelengths of the incident field. However,
in the Born scattering based phase screen model proposed by {\em Wu\/}~(1994),
backscattering of an isolated thin slab is enhanced with respect to 
forward scattering. Although back-scatter is enhanced locally, with 
respect to forward-scatter, it remains an open question whether  
back-scatter is enhanced relative to forward-scatter on a global scale. 
From past computational experience, we know that back scatter is more 
pronounced when the medium changes very rapidly, on the scale length of 
the incident field (which is certainly the case for a thin slab). When 
the individual slabs are bonded to form a thicker region, it is hoped 
that the back scatter will be attenuated through interference.
\item by applying the {\em parabolic or narrow angle\/} approximation in 
the derivation of the forward scattering interaction at each screen, some 
of the gains in accuracy of the wave based theory over the ray based 
theory have been wasted.
\end{itemize}


We now show how to correct the forward scattering model of {\em 
Wu\/}~(1994), for the neglected backscatter using wave field balances   and matrix factorization 
techniques. We also derive single scattering operators, that are local in 
the spatial domain, with out invoking a narrow angle, or {\em 
parabolic\/} approximation. Our approach is based on the {\em 
pseudo-spectral approximation}
\begin{eqnarray} \label{laydSdef}
\lefteqn{
{\mit \delta}{\CB S}^{\rho}(k_y)=
\left({\begin{array}{cc}
{\mit \delta}{\bmi S}_{\uparrow \downarrow }^{\rho}&{\mit \delta}{\bmi S}_{ \uparrow \uparrow }^{\rho}\\ 
{\mit \delta}{\bmi S}_{\downarrow \downarrow }^{\rho}&{\mit \delta}{\bmi S}_{ \downarrow \uparrow}^{\rho}
\end{array}}\right)=
-i\,\left({\begin{array}{cc}
\sqrt{{\bmi E}_{\rho}}&{\bf 0}\\ 
{\bf 0}&\sqrt{{\bmi E}_{\rho}}
\end{array}}\right)\,
\hat{\bmi  D}{}^{T}_{z,\rho}(-{\bf k}_{x},-{k}_{y})
\left({\bmi I}_8\otimes{\bf DFT}^{T}(-{\bf k}_{x})  \right)}\nonumber\\
&&
\times 
\hat{\bmi \Phi}{\mit \Delta}{\bmi C}^{\rho}(k_y) 
\hat{\bmi \Phi}{}^{-1} \left({\bmi I}_8\otimes 
{\bf DFT}({\bf k}_{x}) \right)\hat{\bmi  D}_{z,\rho}({\bf k}_{x},{k}_{y}) \left({\begin{array}{cc}
\sqrt{{\bmi E}_{\rho}}&{\bf 0}\\ 
{\bf 0}&\sqrt{{\bmi E}_{\rho}}
\end{array}}\right)
\end{eqnarray}
of the $6 N_k \times 6 N_k$ array of plane wave reflection/transmission 
transitions from $\mbox{row}_{\rho}, \ 1 \le \rho \le L$ of the 
heterogeneity.
The {\em effective coupling matrix\/} ${\mit \Delta}{\CB C}{}^{\rho}_{eff}(k_y)$ 
for $\mbox{row}_{\rho}$ was defined by (\ref{dCBCrho}). The approximation
(\ref{laydSdef}) neglects the side scattering component of (\ref{laySdef}).




 The net result, is an algorithm
that is competitive, with the methods developed in the earlier 
chapters of this thesis,  particularly for the subclass of problems where 
the heterogeneity is embedded in 
a uniform host, so that the background medium does not support 
reverberation.



 To derive the {\em forward scattering\/} splitting, we re-order the interfacial variables 
${\CB  V}_{{int}}$ in (\ref{schur}) into their $\downarrow\!\!-going$ and 
$\uparrow\!\!-going$ components to get
\begin{equation} \label{updownV}
{\CB  V}_{{int}}\equiv \left[{\CB  V}_{{int, \downarrow}}, {\CB  
V}_{{int, \uparrow}}\right]^{T}
\end{equation}
When we re-arrange (\ref{schur}) consistent with the new ordering 
(\ref{updownV}) of the field variable ${\CB  V}_{{int}}$, we get
\begin{equation} \label{Avb}
\left[{\begin{array}{cc}
{\bmi A}_{ \downarrow \downarrow }&{\bmi A}_{ \downarrow \uparrow }\\ 
{\bmi A}_{ \uparrow \downarrow }&{\bmi A}_{ \uparrow \uparrow }\end{array}}\right]
\left[{\begin{array}{c}{\CB V}_{int, \downarrow }\\ 
{\CB V}_{int, \uparrow }\end{array}}\right]={\bmi b}_{int}
\end{equation} 
where the right hand side vector ${\bmi b}_{int}$ was given previously by 
(\ref{SurfSdef}) for a source in the surface layer, or (\ref{LowerSdef}) for a source in 
the lower halfspace, while the partitions of the coefficient array in 
(\ref{Avb}) are defined as:
\[
{\bmi A}_{ \downarrow \downarrow 
}=\left[{\begin{array}{ccc}\left[{{\bmi I}-{\bmi S}_{ \downarrow \uparrow 
}^{[f,0]} {\mit \delta} {  \bmi S}_{\uparrow \downarrow }^{  
1}}\right]&&\\ -{\bmi S}_{ \downarrow \downarrow }^{1} \left[{\bmi E}_{1}+ {\mit \delta} {  
\bmi S}_{\downarrow\downarrow  }^{  1}\right]&\left[{  {\bmi I}-{\bmi S}_{ 
\downarrow \uparrow }^{1} {\mit \delta} {  \bmi S}_{\uparrow \downarrow 
}^{  2}}\right]&\\ &\ddots  &\ddots\\ &-{\bmi S}_{ \downarrow \downarrow 
}^{L-1} \left[{\bmi E}_{L-1}+{\mit \delta} {  \bmi S}_{ \downarrow \downarrow }^{  L-1}\right]&\left[{  {\bmi I}-{\bmi S}_{ \downarrow \uparrow }^{L-1} {\mit \delta} { 
 \bmi S}_{\uparrow \downarrow }^{  L}}\right]\end{array}}\right];
\] 



\[
{\bmi A}_{ \uparrow \uparrow }=\left[{\begin{array}{ccc}\left[{{\bmi I}-{\bmi S}_{ 
\uparrow \downarrow }^{1} {\mit \delta} {  \bmi S}_{\downarrow \uparrow 
}^{  1}}\right]&-{\bmi S}_{ \uparrow \uparrow }^{1} 
\left[{\bmi E}_{2}+{\mit \delta} {  \bmi S}_{ \uparrow \uparrow }^{ 2}\right]&\\ &\left[{  {\bmi I}-{\bmi S}_{ 
\uparrow \downarrow }^{2} {\mit \delta} {  \bmi S}_{\downarrow \uparrow 
}^{  2}}\right]&  -{\bmi S}_{ \uparrow \uparrow }^{2} 
\left[{\bmi E}_{3}+{\mit \delta} {\bmi S}_{ \uparrow \uparrow }^{3}\right]\\ 
&\ddots &\ddots  \\ 
&&\left[{{\bmi I}-{\bmi S}_{ \uparrow \downarrow }^{[L, \infty   )} 
{\mit \delta} {  \bmi S}_{\downarrow \uparrow }^{  
L}}\right]\end{array}}\right];
\] 

\[
{\bmi A}_{ \uparrow \downarrow }=-\left[{\begin{array}{ccc}{\bmi S}_{ \uparrow 
\downarrow }^{1} 
\left[{\bmi E}_{1}+{\mit \delta} {  \bmi S}_{ \downarrow \downarrow }^{1}\right]&
{  \bmi S}_{\uparrow \uparrow }^{  1}{\mit \delta} {  
\bmi S}_{\uparrow \downarrow }^{  2}&\\ &{  \bmi S}_{\uparrow 
\downarrow }^{  2}\left[{\bmi E}_{2}+{\mit \delta} {  \bmi S}_{ \downarrow 
\downarrow }^{2}\right]&{  \bmi S}_{\uparrow \uparrow }^{  2}{\mit \delta} {  
\bmi S}_{\uparrow \downarrow }^{  3}\\ &\ddots  &\ddots\\ &&{\bmi S}_{ \uparrow 
\downarrow }^{[L, \infty)} \left[{\bmi E}_{L}+{\mit \delta} {  \bmi S}_{ 
\downarrow \downarrow }^{L}\right]\end{array}}\right];
\] 



\[
{\bmi A}_{ \downarrow \uparrow }=-\left[{\begin{array}{ccc}{\bmi S}_{ 
\downarrow \uparrow }^{[f,0]} 
\left[{\bmi E}_{1}+{\mit \delta} {  \bmi S}_{ \uparrow \uparrow }^{1}\right]&&\\ {  \bmi S}_{\downarrow \downarrow }^{  1}{\mit \delta} 
{  \bmi S}_{\downarrow \uparrow }^{  1}&{  \bmi S}_{\downarrow 
\uparrow }^{  1}\left[{\bmi E}_{2}+{\mit \delta} {  \bmi S}_{ \uparrow \uparrow 
}^{2}\right]&\\ &\ddots &\ddots\\ &{ \bmi S}_{\downarrow \downarrow }^{ L-1}{\mit \delta} { 
 \bmi S}_{\downarrow \uparrow }^{  L-1}&{  \bmi S}_{\downarrow 
\uparrow }^{  L-1}\left[{\bmi E}_{L}+{\mit \delta} {  \bmi S}_{ \uparrow 
\uparrow }^{L}\right]\end{array}}\right];
\]
We now split (\ref{Avb}) according to
\begin{equation} \label{Fsplit}
\left[{\begin{array}{cc}{\bmi P}_{ \downarrow \downarrow }^{-1}&{\bf 0}\\ 
{\bmi A}_{ \uparrow \downarrow }&{\bmi P}_{ \uparrow \uparrow 
}^{-1}\end{array}}\right]\left[{\begin{array}{c}{\CB V}_{int, \downarrow 
}\\ {\CB V}_{int, \uparrow 
}\end{array}}\right]={\bmi b}_{int}+\left[{\begin{array}{cc}{\bmi B}_{ \downarrow 
\downarrow }&{\bmi A}_{ \downarrow \uparrow }\\ {\bf 0}&{\bmi B}_{ \uparrow \uparrow 
}\end{array}}\right]\left[{\begin{array}{c}{\CB V}_{int, \downarrow }\\ 
{\CB V}_{int, \uparrow }\end{array}}\right]
\end{equation} 
where:
\[
{\bmi P}_{ \downarrow \downarrow }^{-1}
=\left[{\begin{array}{ccc}{\bmi I}&&\\ -{\bmi S}_{ \downarrow \downarrow }^{1} 
\left[{\bmi E}_{1}+{\mit \delta} {  \bmi S}_{ \downarrow \downarrow }^{1}\right]&
\hspace{-10em} {\bmi I}&\\
\qquad \ddots &  \ddots &\\ 
&-{\bmi S}_{ \downarrow \downarrow 
}^{L-1} \left[{\bmi E}_{L-1}+{\mit \delta} {  \bmi S}_{ \downarrow \downarrow 
}^{L-1}\right]&{\bmi I}\end{array}}\right];
\]
\[
{\bmi P}_{ \uparrow \uparrow }^{-1}
=\left[{\begin{array}{ccc}
{\bmi I}&-{\bmi S}_{ \uparrow \uparrow }^{1} 
\left[{\bmi E}_{2}+{\mit \delta} {  \bmi S}_{ \uparrow \uparrow }^{2}\right]&\\ 
&{\bmi I}&  -{\bmi S}_{ \uparrow \uparrow }^{2} 
\left[{\bmi E}_{3}+{\mit \delta} {  \bmi S}_{ \uparrow \uparrow }^{3}\right]\\ 
&\hspace{8em} \ddots &\ddots  \\ 
&&{\bmi I}\end{array}}\right];
\]
\[
{\bmi B}_{ \downarrow \downarrow }=
\left[{\begin{array}{ccc}{\bmi S}_{ \downarrow \uparrow 
}^{[f,0]} {\mit \delta} {  \bmi S}_{\uparrow \downarrow }^{  
1}&&\\ 
&{\bmi S}_{\downarrow \uparrow }^{1} {\mit \delta} {  \bmi S}_{\uparrow \downarrow 
}^{  2}&\\ 
&\ddots&\\
&&{\bmi S}_{ \downarrow \uparrow }^{L-1} {\mit \delta} { 
 \bmi S}_{\uparrow \downarrow }^{  L}\end{array}}\right];
\] 
\[
{\bmi B}_{ \uparrow \uparrow }
=\left[{\begin{array}{ccc}{\bmi S}_{ 
\uparrow \downarrow }^{1} {\mit \delta} {  \bmi S}_{\downarrow \uparrow 
}^{  1}&&\\ &{\bmi S}_{\uparrow\downarrow }^{2} 
{\mit \delta}{\bmi S}_{\downarrow\uparrow}^{2}&\\
&\ddots&\\ &&{\bmi S}_{ \uparrow \downarrow }^{[L, \infty   )} 
{\mit \delta} {  \bmi S}_{\downarrow \uparrow }^{L}\end{array}}\right];
\]

Multiplying (\ref{Fsplit}) through by the inverse of the left hand side 
array we have 
\begin{equation}\label{LSmatF}
\fbox{$\displaystyle
\left[{\begin{array}{c}{\CB V}_{int, \downarrow }\\ {\CB V}_{int, \uparrow 
}\end{array}}\right]=
\left[{\begin{array}{cc}{\bmi P}_{ \downarrow \downarrow }^{-1}&{\bf 0}\\ 
{\bmi A}_{ \uparrow \downarrow }&{\bmi P}_{ \uparrow \uparrow 
}^{-1}\end{array}}\right]^{-1}{\bmi b}_{int}+
\left[{\begin{array}{cc}{\bmi P}_{ \downarrow \downarrow }^{-1}&{\bf 0}\\ 
{\bmi A}_{ \uparrow \downarrow }&{\bmi P}_{ \uparrow \uparrow 
}^{-1}\end{array}}\right]^{-1}\left[{\begin{array}{cc}{\bmi B}_{ \downarrow 
\downarrow }&{\bmi A}_{ \downarrow \uparrow }\\ {\bf 0}&{\bmi B}_{ \uparrow \uparrow 
}\end{array}}\right]\left[{\begin{array}{c}{\CB V}_{int, \downarrow }\\ 
{\CB V}_{int, \uparrow }\end{array}}\right]
$}
\end{equation} 
which is a {\em forward scattering\/} analogue of the Lippmann-Schwinger 
equation (\ref{LSmat}).  Using the matrix identity
\[
{\left[{\begin{array}{cc}{\bmi P}_{ \downarrow \downarrow }^{-1}&{\bf 0}\\
{\bmi A}_{ \uparrow \downarrow }&{\bmi P}_{ \uparrow \uparrow }^{-1}\end{array}}\right]}^{-1}=
\left[{\begin{array}{cc}{\bmi I}&{\bf 0}\\
{\bf 0}&{\bmi P}_{ \uparrow \uparrow }\end{array}}\right]
\left[{\begin{array}{cc}{\bmi I}&{\bf 0}\\
-{\bmi A}_{ \uparrow \downarrow }&{\bmi I}\end{array}}\right]
\left[{\begin{array}{cc}{\bmi P}_{ \downarrow \downarrow }&{\bf 0}\\{\bf 
0}&{\bmi I}\end{array}}\right]
\]
in (\ref{LSmatF}) we get
\begin{eqnarray} \label{fLSeqn}
\left[{\begin{array}{c}{\CB V}_{int, \downarrow }\\ {\CB V}_{int, \uparrow 
}\end{array}}\right]&=&\left[{\begin{array}{c}{\CB V}_{inc, \downarrow }\\ 
{\CB V}_{inc, \uparrow 
}\end{array}}\right]+\nonumber\\
&&
\left[{\begin{array}{cc}{\bmi I}&{\bf 0}\\
{\bf 0}&{\bmi P}_{ \uparrow \uparrow }\end{array}}\right]
\left[{\begin{array}{cc}{\bmi I}&{\bf 0}\\
-{\bmi A}_{ \uparrow \downarrow }&{\bmi I}\end{array}}\right]
\left[{\begin{array}{cc}{\bmi P}_{ \downarrow \downarrow }&{\bf 0}\\{\bf 
0}&{\bmi I}\end{array}}\right]
\left[{\begin{array}{cc}{\bmi B}_{ \downarrow 
\downarrow }&{\bmi A}_{ \downarrow \uparrow }\\ {\bf 0}&{\bmi B}_{ \uparrow \uparrow 
}\end{array}}\right]\left[{\begin{array}{c}{\CB V}_{int, \downarrow }\\ 
{\CB V}_{int, \uparrow }\end{array}}\right]
\end{eqnarray} 
where
\[
\left[{\begin{array}{c}{\CB V}_{inc, \downarrow }\\ 
{\CB V}_{inc,\uparrow}\end{array}}\right]=
\left[{\begin{array}{cc}{\bmi I}&{\bf 0}\\
{\bf 0}&{\bmi P}_{ \uparrow \uparrow }\end{array}}\right]
\left[{\begin{array}{cc}{\bmi I}&{\bf 0}\\
-{\bmi A}_{ \uparrow \downarrow }&{\bmi I}\end{array}}\right]
\left[{\begin{array}{cc}{\bmi P}_{ \downarrow \downarrow }&{\bf 0}\\{\bf 
0}&{\bmi I}\end{array}}\right]
{\bmi b}_{int}
\] 
is the {\em effective incident field\/}. Due to the sparsity 
pattern of ${\bmi P}_{ \downarrow \downarrow }^{-1}$ and ${\bmi P}_{ 
\uparrow \uparrow }^{-1}$, we can implement the {\em down marching\/} 
operation ${\bmi P}_{ \downarrow \downarrow }$ and the {\em up marching\/} 
operation ${\bmi P}_{ \uparrow \uparrow }$ without any matrix inversions.


\section{A computable representation of the elastic wave Lippmann-Schwinger equation}

The solution of the matrix Lippmann-Schwinger equation (\ref{LSmat}), may be 
written as
\begin{equation} \label{MSVint}
{\CB  V}_{{int}}
=\left[{\bmi  I} - {\CB Q}^{\partial} {\mit \Delta} {\CB  S}\right]^{-1}
{\CB  V}_{{inc}}
={\CB  V}_{{inc}}+
{\CB Q}^{\partial} {\mit \Delta} {\CB  S}\left[{\bmi I} - 
{\CB Q}^{\partial} {\mit \Delta} {\CB  S}\right]^{-1}
{\CB  V}_{{inc}}
\end{equation}
where the incident field ${\CB  V}_{{inc}}$  was given by (\ref{sourceV}),
the planewave propagator for a stratified halfspace ${\CB Q}^{\partial}$ 
was given by (\ref{PstratDef}) and finally, the block diagonal plane wave 
coupling array ${\mit \Delta} {\CB  S}$ was given in (\ref{GSdef}).
From (\ref{GSdef1}) we have
\begin{equation} \label{GSdef2}
{\mit \Delta}{\CB  S}(k_y) \ \equiv\  -i\, {\bmi \Theta}^{z}_{ou} 
{\mit \Delta}{\CB T}(k_y) {\bmi \Theta}^{z}_{in}
\end{equation}
where ${\bmi \Theta}^{z}_{ou}$, 
${\mit \Delta}{\CB T}(k_y)$ and  ${\bmi \Theta}^{z}_{in}$ were defined by 
(\ref{Thetaouz})--(\ref{Thetainz}) of {\sc Box } {\bf 5.2}.

Using (\ref{GSdef2}) for ${\mit \Delta}{\CB  S}(k_y)$ in (\ref{MSVint}), we get
\begin{eqnarray} \label{MSVint1}
{\CB  V}_{{int}}
&=&{\CB  V}_{{inc}}-i\,
{\CB Q}^{\partial}  {\bmi \Theta}^{z}_{ou} {\mit \Delta}{\CB T}
{\bmi \Theta}^{z}_{in}\left[{\bmi I} + i\, {\CB Q}^{\partial} 
{\bmi \Theta}^{z}_{ou} {\mit \Delta}{\CB T}
{\bmi \Theta}^{z}_{in}\right]^{-1}
{\CB  V}_{{inc}}\nonumber\\
&=&{\CB  V}_{{inc}}-i\,
{\CB Q}^{\partial}  {\bmi \Theta}^{z}_{ou} {\mit \Delta}{\CB T}
\left[{\bmi I} + i\, {\bmi \Theta}^{z}_{in} {\CB Q}^{\partial} 
{\bmi \Theta}^{z}_{ou} {\mit \Delta}{\CB T}
\right]^{-1}
{\bmi \Theta}^{z}_{in}{\CB  V}_{{inc}}
\end{eqnarray}



Using (\ref{layerT}) 
for the diagonal blocks 
 in ${\mit \Delta}{\CB T}$, (\ref{CBT}) may 
be written as
\begin{equation} \label{CBT1}
{\mit \Delta}{\CB T}(k_y)= 
{\mit \Delta}{\CB C}{}_{eff}(k_y)\left[{\bmi I}- 
\rf{\CB P}\!{}^{x}_{r,s}\cdot {\mit \Delta}{\CB 
C}{}_{eff}(k_y)  \right]^{-1}
\end{equation}
where ${\mit \Delta}{\CB C}{}_{eff}(k_y)$ is the coupling array for the 
whole layer stack, that was defined in (\ref{dCBC}), and 
$\rf{\CB P}\!{}^{x}_{r,s}$ is the block diagonal array
\begin{equation} \label{CBPxaugL}
\rf{\CB P}\!{}^{x}_{r,s}=\mbox{diag}\left[
\rf{\CB P}\!{}^{x}_{r,s, 1}, \cdots,
\rf{\CB P}\!{}^{x}_{r,s, \rho}, \cdots,
\rf{\CB P}\!{}^{x}_{r,s, L}
\right]
\end{equation}
comprised from the sideways propagators $\rf{\CB P}\!{}^{x}_{r,s, \rho}$ 
of (\ref{CBPxaug}), for each of the $L$-layers of the 
reference stratification.



Using (\ref{CBT1}) for ${\mit \Delta}{\CB T}(k_y)$ in (\ref{MSVint1}) gives
\begin{eqnarray} \label{MSVint2}
\lefteqn{{\CB  V}_{{int}}
={\CB V}_{{inc}}-i\, {\CB Q}^{\partial} {\bmi \Theta}^{z}_{ou} 
{\mit \Delta}{\CB C}{}_{eff}(k_y) \left[{\bmi I}- \rf{\CB P}\!{}^{x}_{r,s}\cdot {\mit 
\Delta}{\CB C}{}_{eff}(k_y) \right]^{-1}}\nonumber\\ 
&&\times \bigg({\bmi I} + i\, 
 {\bmi \Theta}^{z}_{in} {\CB Q}^{\partial} 
{\bmi \Theta}^{z}_{ou} 
{\mit \Delta}{\CB C}{}_{eff}(k_y) \left[{\bmi I}- \rf{\CB P}\!{}^{x}_{r,s}\cdot {\mit \Delta}{\CB 
C}{}_{eff}(k_y) \right]^{-1} \bigg)^{-1} {\bmi \Theta}^{z}_{in}{\CB 
V}_{{inc}} 
\end{eqnarray}

We introduce the up/down propagator
\begin{equation} \label{rfCBPzrs}
\rf{\CB P}\!{}^{z}_{r,s}\equiv
-i\,{\bmi \Theta}^{z}_{in} {\CB Q}^{\partial} 
{\bmi \Theta}^{z}_{ou}
\end{equation}
where ${\bmi \Theta}^{z}_{ou}$, 
and  ${\bmi \Theta}^{z}_{in}$ were defined by 
(\ref{Thetaouz}) and (\ref{Thetainz}) respectively,  of {\sc Box } {\bf 5.2}.
Using (\ref{rfCBPzrs}) for $\rf{\CB P}\!{}^{z}_{r,s}$ in
 (\ref{MSVint2})  we get
 \begin{samepage}
\begin{eqnarray} \label{MSVint30}
{\CB  V}_{{int}}&=&{\CB  V}_{{inc}}-i\,
{\CB Q}^{\partial}  {\bmi \Theta}^{z}_{ou} 
{\mit \Delta}{\CB C}{}_{eff}(k_y)\nonumber\\
&&\times
\bigg({\bmi I} -  
\left[\rf{\CB P}\!{}^{z}_{r,s}+\rf{\CB P}\!{}^{x}_{r,s}\right]
{\mit \Delta}{\CB C}{}_{eff}(k_y)\bigg)^{-1} 
{\bmi \Theta}^{z}_{in}{\CB  V}_{{inc}}
\end{eqnarray}
\end{samepage}



We assume that both the source and receiver positions 
are excluded from the depth interval $z_0 \le z 
\le z_L$ of the heterogeneity, so that (\ref{MSVint30}) may be written as
\begin{eqnarray} \label{MSVint3}
{\CB  V}_{{int},R}&=&{\CB  V}_{{inc},R}-i\,
{\CB Q}^{\partial}_{R}  {\bmi \Theta}^{z}_{ou} 
{\mit \Delta}{\CB C}{}_{eff}(k_y)\nonumber\\
&&\times
\bigg({\bmi I} -  
\left[\rf{\CB P}\!{}^{z}_{r,s}+\rf{\CB P}\!{}^{x}_{r,s}\right]
{\mit \Delta}{\CB C}{}_{eff}(k_y)\bigg)^{-1} 
{\bmi \Theta}^{z}_{in}{\CB  V}_{{inc}}
\end{eqnarray}
where ${\CB  V}_{{int},R}$ is the vector of plane wave amplitudes in the 
receiver layer, ${\CB  V}_{{inc},R}$ is the corresponding vector when there 
is no superimposed heterogeneity and ${\CB Q}^{\partial}_{R}$ is the plane 
wave propagator mapping the secondary source radiation induced by the 
superimposed lateral heterogeneity 
from the depth interval $z_0 \le z 
\le z_L$ to the receiver layer. More explicitly,  ${\CB 
Q}^{\partial}_{R}$ may be written as
\[
{\CB Q}^{\partial}_{R} \equiv \left[\begin{array}{ccccc} 
{\CB Q}^{\partial}_{R\leftarrow 1}&
,\cdots,&{\CB Q}^{\partial}_{R\leftarrow l}&,\cdots,&
{\CB Q}^{\partial}_{R\leftarrow L}  \end{array}\right]
\]
where the partition ${\CB Q}^{\partial}_{R\leftarrow l}$ 
appropriate to layer-$l$ is defined as
\begin{eqnarray} \label{CBQldef}
{\CB Q}^{\partial}_{R\leftarrow l}&=&
\left[{\begin{array}{c}
{\bmi S}{}_{\downarrow\uparrow }^{   (f,{R-1}+)}{{\bmi E}}_{R}\\ 
{\bmi  I}_{ 3 N_k}\end{array}}\right]   
{\left({\bmi  I}_{ 3 N_k}-
{\bmi S}{}_{\uparrow\downarrow }^{ (R-,L+)} {\bmi S}{}_{\downarrow\uparrow 
}^{(f,R-)}\right)}^{-1} {\bmi S}{}_{\uparrow\uparrow }^{ (R-,{l-1}+)}\nonumber\\
&&\times
{\left({{\bmi  I}_{
3 N_k}-{\bmi S}{}_{\uparrow\downarrow }^{ ({l-1}+,L+)}
{{\bmi S}}{}_{\downarrow\uparrow }^{ (R-,{l-1}+)}}\right)}^{-1}\left[{\begin{array}{cc}
{\bmi  I}_{ 3 N_k}&\quad{{\bmi E}}_{l}{\bmi 
S}{}_{\uparrow\downarrow }^{(l-, 
L+)} \end{array}}\right]
\end{eqnarray}
for a receiver at $(z_R,x_R)$ in the middle of layer-$R$, 
which is above to the zone of heterogeneity $z_0 \le z \le z_L$.

 We synthesize the displacement and traction at
the receiver location from the plane waves ${\CB  V}_{{int},R}$ at that
 depth by operating on both sides of (\ref{MSVint3}) with the 
 displacement/traction  synthesis operator
\begin{equation} \label{ThetazR}
{\bmi \Psi}^{z}_{ou,R}(x_R)\equiv \left({\bmi I}_6\otimes 
{\bmi F}(x_{R};{\bf k}_{x})\right){\bmi  D}_{z,R}({\bf k}_{x},{k}_{y})
\left({\begin{array}{cc}
\sqrt{{\bmi E}_{R}}&{\bf 0}\\ 
{\bf 0}&\sqrt{{\bmi E}_{R}}
\end{array}}\right)
\end{equation}
where ${\bmi E}_{R}$ is the phase shift array (\ref{Ejdef}) appropriate to the 
receiver layer,
  ${{\bmi F}}$ was defined in (\ref{Fdef}) and ${\bmi  D}_{z,R}$ is 
(\ref{Dsuperdef}) evaluating using the reference material properties appropriate to 
the receiver layer.



Operating on both sides of (\ref{MSVint3}) with ${\bmi 
\Psi}^{z}_{in,R}$ from (\ref{ThetazR}) gives
\begin{equation} \label{LSstrat}
\fbox{$\displaystyle
{\CB G}=\rf{\CB G}+
\rf{\CB P}\!{}^{z,R}_{s} 
{\mit \Delta}{\CB C}{}_{eff}(k_y)
\bigg({\bmi I} -  
\left[\rf{\CB P}\!{}^{z}_{r,s}+\rf{\CB P}\!{}^{x}_{r,s}\right]
{\mit \Delta}{\CB C}{}_{eff}(k_y)\bigg)^{-1} 
\rf{\CB G}\!{}^{z}_{r}
$}
\end{equation}
We have defined $\rf{\CB G}\!{}^{z}_{r}$ in (\ref{LSstrat}) as
\[
\rf{\CB G}\!{}^{z}_{r} \equiv {\bmi \Theta}^{z}_{in}{\CB  
V}_{{inc}}\equiv
\left[\begin{array}{c}
\rf{\CB G}\!{}^{z,1}_{r}\\
\vdots\\
\rf{\CB G}\!{}^{z,l}_{r}\\
\vdots\\
\rf{\CB G}\!{}^{z,L}_{r}
\end{array}  \right]
\]
with the partition $\rf{\CB G}\!{}^{z,l}_{r}$ for  layer-$l$ defined as
\begin{eqnarray} \label{Gincdown}
\lefteqn{
\rf{\CB G}\!{}^{z,l}_{r}=\hat{\bmi \Phi}{}^{-1}\left({\bmi I}_8\otimes 
{\bf DFT}({\bf k}_{x}) \right)\hat{\bmi  D}_{z,l}({\bf k}_{x},{k}_{y})
\left[{\begin{array}{c}\sqrt{{\bmi E}_{l}}\\ 
\sqrt{{\bmi E}_{l}} {\bmi S}{}_{\uparrow\downarrow }^{[l, 
\infty)}{{\bmi E}}_{l}\end{array}}\right] 
}\\ 
&&
\times
{\left({\ {\bmi  I}_{ 3 N_k} -
{\bmi S}{}_{\downarrow\uparrow }^{({s} , {l-1} ]}{\bmi S}{}_{\uparrow 
\downarrow
}^{( {l-1} , \infty )} }\right)}^{-1} {\bmi S}{}_{\downarrow
\downarrow }^{( {s} , {l-1} ]}{\left({ {\bmi  I}_{ 3 N_k} - 
{\bmi S}{}_{\downarrow \uparrow }^{[ f , {s} )} {\bmi S}{}_{\uparrow
\downarrow }^{( {s} , \infty )} }\right)}^{-1}
\left[{\bmi S}{}_{\downarrow \uparrow }^{[f, {s})}{\bmi \Sigma}_{\uparrow}+ {\bmi 
\Sigma}_{\downarrow} \right]\nonumber
\end{eqnarray}
when the source is in the surface zone, or as
\begin{eqnarray} \label{Gincup}
\lefteqn{
\rf{\CB G}\!{}^{z,l}_{r}=
\hat{\bmi \Phi}{}^{-1}\left({\bmi I}_8\otimes 
{\bf DFT}({\bf k}_{x}) \right)\hat{\bmi  D}_{z,l}({\bf k}_{x},{k}_{y})
\left[{\begin{array}{c}
\sqrt{{\bmi E}_{l}}{\bmi S}{}_{\downarrow\uparrow }^{   [f,{l-1}]}{{\bmi E}}_{l}\\ 
\sqrt{{\bmi E}_{l}}\end{array}}\right] 
}\nonumber\\
&&
\times
{\left({\bmi  I}_{ 3 N_k}-
{\bmi S}{}_{\uparrow\downarrow }^{ [l,{s})} {\bmi S}{}_{\downarrow\uparrow 
}^{ 
[f,l)}\right)}^{-1}\
 {\bmi S}{}_{\uparrow\uparrow }^{ [l,z_{s})}
{\left({{\bmi  I}_{3 N_k}-{\bmi S}{}_{\uparrow\downarrow }^{ ({s},\infty)}
{{\bmi S}}{}_{\downarrow\uparrow }^{ [f,{s})}}\right)}^{-1}
\left[ {\bmi \Sigma}_{\uparrow} +
{\bmi S}{}_{\uparrow\downarrow }^{({s}, \infty)} {\bmi \Sigma}_{\downarrow}
\right].
\end{eqnarray}
when the source is in the lower halfspace.




The $\rf{\CB P}\!{}^{z,R}_{s}$ array in (\ref{LSstrat})  is defined as
\begin{equation} \label{rfCBPzRs}
\rf{\CB P}\!{}^{z,R}_{s} \equiv -i\,{\bmi \Psi}^{z}_{in,R}
{\CB Q}^{\partial}_{R}  {\bmi \Theta}^{z}_{ou}
\equiv\left[\begin{array}{ccccc} \left(\rf{\CB P}\!{}^{z,R}_{s}\right)_{1}&
,\cdots,&\left(\rf{\CB P}\!{}^{z,R}_{s}\right)_{l}&,\cdots,&
\left(\rf{\CB P}\!{}^{z,R}_{s}\right)_{L}  \end{array}\right]
\end{equation}
where the partition $\left(\rf{\CB P}\!{}^{z,R}_{s}\right)_{l}$ 
appropriate to layer-$l$ is defined as
\begin{eqnarray} \label{rfCBPzRsl}
\lefteqn{
\left(\rf{\CB P}\!{}^{z,R}_{s}\right)_{l}=-i\,
 \left({\bmi I}_6\otimes 
{\bmi F}(x_{R};{\bf k}_{x})\right){\bmi  D}_{z,R}({\bf k}_{x},{k}_{y})
\left[{\begin{array}{c}
\sqrt{{\bmi E}_{R}}{\bmi S}{}_{\downarrow\uparrow }^{   (f,{R-1}+)}{{\bmi E}}_{R}\\ 
\sqrt{{\bmi E}_{R}}\end{array}}\right]   
 }\nonumber\\
&&
\times{\left({\bmi  I}_{ 3 N_k}-
{\bmi S}{}_{\uparrow\downarrow }^{ (R-,L+)} {\bmi S}{}_{\downarrow\uparrow 
}^{(f,R-)}\right)}^{-1}\
 {\bmi S}{}_{\uparrow\uparrow }^{ (R-,{l-1}+)}
{\left({{\bmi  I}_{
3 N_k}-{\bmi S}{}_{\uparrow\downarrow }^{ ({l-1}+,L+)}
{{\bmi S}}{}_{\downarrow\uparrow }^{ (R-,{l-1}+)}}\right)}^{-1}\nonumber\\
&&
\times\left[{\begin{array}{cc}
\sqrt{{\bmi E}_{l}}&\quad{{\bmi E}}_{l}{\bmi 
S}{}_{\uparrow\downarrow }^{(l-, 
L+)} \sqrt{{\bmi E}_{l}}\end{array}}\right]
\,\hat{\bmi  D}{}^{T}_{z,l}(-{\bf k}_{x},-{k}_{y})\left({\bmi I}_8\otimes
{\bf DFT}^{T}(-{\bf k}_{x})  \right)\hat{\bmi \Phi}
\end{eqnarray}
for a receiver in the surface layers.

The array $\rf{\CB G}$ in (\ref{LSstrat}) is the displacement-traction 
Green's tensor for the plane stratified reference medium. To find the 
corresponding Green's tensor ${\CB G}$ for the heterogeneous medium, we 
must evaluate the {\em scattering transition\/}
\[
\rf{\CB P}\!{}^{z,R}_{s} 
{\mit \Delta}{\CB C}{}_{eff}(k_y)
\bigg({\bmi I} -  
\left[\rf{\CB P}\!{}^{z}_{r,s}+\rf{\CB P}\!{}^{x}_{r,s}\right]
{\mit \Delta}{\CB C}{}_{eff}(k_y)\bigg)^{-1} 
\rf{\CB G}\!{}^{z}_{r}
\]
in (\ref{LSstrat}). This task is considered at length in chapter 
\ref{variational}. 

The key step in the solution algorithms of chapter \ref{variational} 
entails the evaluation of the matrix/vector product
\[
\bigg({\bmi I} -  
\rf{\CB P}_{r,s}
{\mit \Delta}{\CB C}{}_{eff}(k_y)\bigg) \cdot \hat{\CB B}
\]
for an arbitrary $8 L N_x-$vector  $\hat{\CB B}$, where the propagator 
$\rf{\CB P}_{r,s}$ is defined as 
\begin{eqnarray} \label{CBPdef0}
\rf{\CB P}_{r,s}&=&\rf{\CB P}\!{}^{z}_{r,s}+\rf{\CB P}\!{}^{x}_{r,s}\nonumber\\
&=&-i\,{\bmi \Theta}^{z}_{in} {\CB Q}^{\partial} 
{\bmi \Theta}^{z}_{ou}
+\mbox{diag}\left[
\rf{\CB P}\!{}^{x}_{r,s, 1}, \cdots,
\rf{\CB P}\!{}^{x}_{r,s, \rho}, \cdots,
\rf{\CB P}\!{}^{x}_{r,s, L}
\right]
\end{eqnarray}
where we have used the definitions of $\rf{\CB P}\!{}^{z}_{r,s}$ and 
$\rf{\CB P}\!{}^{x}_{r,s}$ from (\ref{rfCBPzrs}) and (\ref{CBPxaugL}) 
respectively. 
The block diagonal array ${\mit \Delta}{\CB 
C}{}_{eff}(k_y)$ from (\ref{dCBC}),  is comprised from the coupling screens of the 
individual layers of the heterogeneity. Due to the block diagonal nature of
${\mit \Delta}{\CB C}{}_{eff}(k_y)$ the evaluation of 
\[
\hat{\CB F}={\mit \Delta}{\CB C}{}_{eff}(k_y)\cdot \hat{\CB B}
\]
is both straightforward to code, and efficient to implement.

In algorithm {\bf 5.2} a pseudo-code for the implementation of $\rf{\CB 
P}\!{}^{x}_{r,s} \hat{\CB F}$ is presented, while the pseudo-code for 
$\rf{\CB P}\!{}^{z}_{r,s} \hat{\CB F}$ is given in algorithm {\bf 5.3}, 
where $\hat{\CB F}={\mit \Delta}{\CB C}{}_{eff}(k_y)\cdot \hat{\CB B}$.


%%%%%%%%%%%%%%%%%%%%%%%%%%%%%%%%%%%%%%%%%%%%%%%%%%%%%%%%%%%%%%%%%%%%%%%%%%%%%%%%%%%%%%%%%%%
%%%%%%%%%%%%%%%%%%%%%%%%%%%%%%%%%%%%%%%%%%%%%%%%%%%%%%%%%%%%%%%%%%%%%%%%%%%%%%%%%%%%%%%%%%%
%%%%%%%%%%%%%%%%%%%%%%%%%%%%%%%%%%%%%%%%%%%%%%%%%%%%%%%%%%%%%%%%%%%%%%%%%%%%%%%%%%%%%%%%%%%
\begin{figure}
%\vspace*{10em}

\fbox{
\begin{minipage} {140.0mm}
{\bf  Algorithm}~{\bf 5.2} {\sc Implement the sideways sweeps 
$\hat{\CB B}=\rf{\CB P}\!{}^{x}_{r,s}\hat{\CB F}$ within each row of the 
heterogeneity lattice}\newline
\vspace{4mm}



Given that:
\begin{itemize}
\item the $3 \bar{N}_k \times 3 \bar{N}_k$
diagonal array ${\bf E}_{x,\rho}$, is the horizontal phase shift array ${\bf 
E}_{x}$ of (\ref{Exdef})

\item the planewave propagator ${\CB Q}{}^{x}_{\rho \leftarrow \rho}$
is given by (\ref{CBQxdef})

\item the $8 \bar{N}_k \times 6 \bar{N}_k$ array $\hat{\bmi  
D}_{x,\rho}({\bf k}_{z},{k}_{y})$ was given by (\ref{hatDxsuperdef})

\item the $\bar{N}_k-$element row vector $\overline{\bmi F}(0)$ was given 
by (\ref{Fbardef})

\end{itemize}
The appended subscript $\rho$ indicates 
that the array should be constructed using the reference 
material parameters appropriate for layer-$\rho$.

We evaluate
\[
\hat{\CB B}=
\mbox{diag}\left[
\rf{\CB P}\!{}^{x}_{r,s, 1}, \cdots,
\rf{\CB P}\!{}^{x}_{r,s, \rho}, \cdots,
\rf{\CB P}\!{}^{x}_{r,s, L}
\right]\hat{\CB F}
\]
where $\hat{\CB F}$ is an arbitrary $8 L N_x-$vector,
$\rf{\CB P}\!{}^{x}_{r,s, \rho}$ was  defined in (\ref{CBPxaug}) as
\begin{eqnarray*}
\!\!\!\!\!\!\!\!\!
\lefteqn{\rf{\CB P}\!{}^{x}_{r,s,\rho} =}\nonumber\\
&&\!\!\!\!\!\!\!\!\!
 -{i} {\bmi I}_{N_x} \otimes
\left\{\left({\bmi I}_{8}\otimes \overline{\bmi F}(0)  \right)
 \hat{\bmi  D}_{x,\rho}({\bf k}_{z},{k}_{y})\right\}
 {\CB Q}{}^{x}_{\rho \leftarrow \rho}{\bmi I}_{N_x} \otimes 
 \left\{\hat{\bmi  D}{}_{x,\rho}^{T}(-{\bf k}_{z},-{k}_{y})
\,\left({\bmi I}_{8}\otimes \overline{\bmi F}{}^{T}(0)  \right)
\right\}
\end{eqnarray*}


\begin{itemize}
\item  (*{\em downsweep\/ }*) 
\item {\bf  For} $\rho=1,\cdots,L$
\begin{itemize}
\item {\bf  For} $\chi=1,\cdots,N_x$
\begin{itemize}
\item $
\left[\begin{array}{c}
{\bmi \Sigma}^{(\rho,\chi)}_{\leftarrow}\\
{\bmi \Sigma}^{(\rho,\chi)}_{\rightarrow}
\end{array}\right]=-i\,
\hat{\bmi  D}{}_{x,\rho}^{T}(-{\bf k}_{z},-{k}_{y})
\,\left({\bmi I}_8\otimes \overline{\bmi F}{}^{T}(0)  \right)
\hat{\bmi f}_{N_x(\rho-1)+\chi}
$
\end{itemize}
\item {\bf  End} !$\chi-$loop
\end{itemize}

\item (* {\em right sweep\/} *)

\begin{itemize}
\item $
{\mit \Delta} {\bmi  V}_{\rightarrow }^{  
(\rho,1)}={\bf 0}$

\item {\bf  For} $\chi=2,\cdots,N_x$
\begin{itemize}
\item $
{\mit \Delta} {\bmi  V}_{\rightarrow }^{(\rho,\chi)} ={\bf E}_{x,\rho}\left({{ \bmi
\Sigma}_{ \rightarrow }^{(\rho,\chi-1)}+ {\mit \Delta} { \bmi V}_{\rightarrow 
}^{(\rho,\chi-1)}}
\right)
$
\end{itemize}
\item {\bf  End} !$\chi-$loop
\end{itemize}

\item (* {\em left sweep\/} *)
\begin{itemize}
\item $
{\mit \Delta} {\bmi  V}_{\leftarrow }^{  
(\rho,N_x)}={\bf 0}$

\item {\bf  For} $\chi=Nx-1,\cdots,1$
\begin{itemize}
\item $
{\mit \Delta} { \bmi V}_{\leftarrow }^{(\rho,\chi)}  = {\bf E}_{x,\rho}\left({{ \bmi
\Sigma}_{ \leftarrow }^{(\rho,\chi+1)}+ {\mit \Delta} {\bmi  V}_{\leftarrow }^{  
(\rho,\chi+1)}}\right)
$
\end{itemize}
\item {\bf  End} !$\chi-$loop
\end{itemize}

\item (* {\em right sweep\/} *)

\begin{itemize}
\item {\bf  For} $\chi=1,\cdots,N_x$
\begin{itemize}
\item $
{\bmi b}^{(\rho,\chi)}=
\left({\bmi I}_8\otimes \overline{\bmi F}(0)  \right)
 \hat{\bmi  D}_{x,\rho}({\bf k}_{z},{k}_{y})
 \left[\begin{array}{c}
{\mit \Delta}{\bmi V}^{(\rho,\chi)}_{\rightarrow}\\
{\mit \Delta}{\bmi V}^{(\rho,\chi)}_{\leftarrow}
\end{array}\right]
$

\end{itemize}
\item {\bf  End} !$\chi-$loop
\end{itemize}

\item {\bf  End} !$\rho-$loop

\end{itemize}

\vspace{2 mm}
\end{minipage}
}
\end{figure}
%%%%%%%%%%%%%%%%%%%%%%%%%%%%%%%%%%%%%%%%%%%%%%%%%%%%%%%%%%%%%%%%%%%%%%%%%%%%%%%%%%%%%%%%%%%
%%%%%%%%%%%%%%%%%%%%%%%%%%%%%%%%%%%%%%%%%%%%%%%%%%%%%%%%%%%%%%%%%%%%%%%%%%%%%%%%%%%%%%%%%%%
%%%%%%%%%%%%%%%%%%%%%%%%%%%%%%%%%%%%%%%%%%%%%%%%%%%%%%%%%%%%%%%%%%%%%%%%%%%%%%%%%%%%%%%%%%%
%%%%%%%%%%%%%%%%%%%%%%%%%%%%%%%%%%%%%%%%%%%%%%%%%%%%%%%%%%%%%%%%%%%%%%%%%%%%%%%%%%%%%%%%%
%%%%%%%%%%%%%%%%%%%%%%%%%%%%%%%%%%%%%%%%%%%%%%%%%%%%%%%%%%%%%%%%%%%%%%%%%%%%%%%%%%%%%%%%%%%
%%%%%%%%%%%%%%%%%%%%%%%%%%%%%%%%%%%%%%%%%%%%%%%%%%%%%%%%%%%%%%%%%%%%%%%%%%%%%%%%%%%%%%%%%%%
%%%%%%%%%%%%%%%%%%%%%%%%%%%%%%%%%%%%%%%%%%%%%%%%%%%%%%%%%%%%%%%%%%%%%%%%%%%%%%%%%%%%%%%%%%%

\begin{figure}
%\vspace*{10em}

\fbox{
\begin{minipage} {140.0mm} {\bf  Algorithm}~{\bf 5.3}
{\sc Implement the up/down sweeps $\hat{\CB B}=\rf{\CB P}\!{}^{z}_{r,s}\hat{\CB F}$
between rows of the heterogeneity lattice}\newline
\vspace{4mm}


Given:
\begin{itemize}
\item the $3 {N}_k \times 3 {N}_k$ diagonal, vertical phase shift arrays  
${\bmi  E}_{\rho}$ , for each $1 \le \rho \le L$, that were defined in 
(\ref{Ejdef})

\item the $8 {N}_k \times 6 {N}_k$ arrays $\hat{\bmi  
D}_{z,\rho}({\bf k}_{z},{k}_{y})$, for each $1 \le \rho \le L$,
 that were given by (\ref{hatDsuperdef})

\item the $N_k \times N_k$, DFT array ${\bf DFT}$ that was given 
by (\ref{DFTdef})

\item the upsampling and downsampling operators $\hat{\bmi 
\Phi}\left[\cdot \right]$ and $\hat{\bmi \Phi}{}^{-1}\left[\cdot\right]$ respectively, 
that are implemented using Alg.{\bf 4.1}
 
\item the interfacial reflection/transmission arrays
\[
\left[{\begin{array}{cc}{\bmi S}^{\rho}_{ 
\uparrow \downarrow }&{\bmi S}^{\rho}_{ \uparrow \uparrow }\\ {\bmi S}^{\rho}_{ \downarrow 
\downarrow }&{\bmi S}^{\rho}_{ \downarrow \uparrow 
}\end{array}}\right]
\]
constructed in {\sc Box} {\bf 5.1}, for each interface-$\rho$ with $1 \le \rho \le L-1$.

\item the net downward incident  reflection response ${\bmi S}_{ \uparrow\downarrow 
}^{[L,\infty)}$ of the lower halfspace

\item the net upward incident reflection response ${\bmi S}_{ \downarrow \uparrow
}^{[f,0]}$ of the free surface zone

\end{itemize}
we evaluate
\[
\hat{\CB B}\equiv \rf{\CB P}\!{}^{z}_{r,s}\cdot\hat{\CB F}\equiv
-i\,{\bmi \Theta}^{z}_{in} {\CB Q}^{\partial} 
{\bmi \Theta}^{z}_{ou}\cdot\hat{\CB F}
\]
where:
\begin{itemize}

\item $\hat{\CB F}$ is an arbitrary $8 L N_x-$vector with the $8 
N_x-$subvector components $\hat{\CB F}_{\rho}$ for $\mbox{row}_{\rho}$ of 
heterogeneity, for each   $1 \le \rho \le L$.
In turn, $\hat{\CB F}_{\rho}$ is  comprised from the $8-$vectors, 
${\bmi f}_{N_x(\rho-1)+\chi}$ for each  $1  \le \chi \le N_x$
 
\item the plane wave propagator ${\CB Q}^{\partial}$ was given by (\ref{PstratDef}) as
\[
{\CB Q}^{\partial}\equiv 
({\bmi  I}-{\CB  S}_{{int}}{\CB  E})^{-1}{\CB  S}_{{int}};
\]

\item ${\bmi \Theta}^{z}_{ou}$ was defined in (\ref{Thetaouz}) as
\begin{eqnarray*}
\lefteqn{{\bmi \Theta}^{z}_{ou}\equiv
\mbox{diag}\left[
\left({\begin{array}{cc}
\sqrt{{\bmi E}_{1}}&{\bf 0}\\ 
{\bf 0}&\sqrt{{\bmi E}_{1}}
\end{array}}\right)\,\hat{\bmi  D}{}^{T}_{z,1}(-{\bf k}_{x},-{k}_{y})\left({\bmi I}_8\otimes
{\bf DFT}^{T}(-{\bf k}_{x})  \right)\hat{\bmi \Phi} ,\right.}\\
&& \hspace{6em}
\left. \cdots,
\left({\begin{array}{cc}
\sqrt{{\bmi E}_{L}}&{\bf 0}\\ 
{\bf 0}&\sqrt{{\bmi E}_{L}}
\end{array}}\right)\,\hat{\bmi  D}{}^{T}_{z,L}(-{\bf k}_{x},-{k}_{y})\left({\bmi I}_8\otimes
{\bf DFT}^{T}(-{\bf k}_{x})  \right)\hat{\bmi \Phi}   
\right];
\end{eqnarray*}

\item ${\bmi \Theta}^{z}_{in}$ was defined in (\ref{Thetainz}) as
\begin{eqnarray*}
\lefteqn{{\bmi \Theta}^{z}_{in}\equiv
\mbox{diag}\left[
\hat{\bmi \Phi}{}^{-1}\left({\bmi I}_8\otimes 
{\bf DFT}({\bf k}_{x}) \right)\hat{\bmi  D}_{z,1}({\bf k}_{x},{k}_{y})
\left({\begin{array}{cc}
\sqrt{{\bmi E}_{1}}&{\bf 0}\\ 
{\bf 0}&\sqrt{{\bmi E}_{1}}
\end{array}}\right),\right.}\\
&&\hspace{9em}
\left. \cdots,
\hat{\bmi \Phi}{}^{-1}\left({\bmi I}_8\otimes 
{\bf DFT}({\bf k}_{x}) \right)\hat{\bmi  D}_{z,L}({\bf k}_{x},{k}_{y})
\left({\begin{array}{cc}
\sqrt{{\bmi E}_{L}}&{\bf 0}\\ 
{\bf 0}&\sqrt{{\bmi E}_{L}}
\end{array}}\right)
\right].
\end{eqnarray*}

\end{itemize}

\begin{itemize}
\item{\bf Begin Algorithm}
\item  (*{\em bottom of layer stack\/ }*) 


\item ${\bmi R}_{ \uparrow \downarrow }^{L}\equiv 
{\bmi S}_{ \uparrow\downarrow }^{[L,\infty)}$

\item $
\left[{\begin{array}{c}{\bmi \Sigma}^{L}_{\uparrow}\\ 
{\bmi \Sigma}^{L}_{\downarrow}
\end{array}}\right]=-i\,\left({\begin{array}{cc}
\sqrt{{\bmi E}_{L}}&{\bf 0}\\ 
{\bf 0}&\sqrt{{\bmi E}_{L}}
\end{array}}\right)\,\hat{\bmi  D}{}^{T}_{z,L}(-{\bf k}_{x},-{k}_{y})\left({\bmi I}_8\otimes
{\bf DFT}^{T}(-{\bf k}_{x})  \right)\hat{\bmi \Phi}\left[{\CB F}_{L}  \right]
$

\item  $
{\bmi t}_{ \uparrow }^{L}=
{\bmi R}_{ \uparrow \downarrow }^{L}
{\bmi \Sigma}^{L}_{\downarrow}$


\end{itemize}

\vspace{2mm}

\end{minipage}
}
\end{figure}

\begin{figure}
%\vspace*{10em}
\fbox{
\begin{minipage} {140.0mm}{\bf  Algorithm}~{\bf 5.3} {\em 
continued\/}\newline

\begin{itemize}


\begin{itemize}
\item  (*{\em upsweep\/ }*) 
\end{itemize}
\item {\bf  For} $\rho=L-1,\cdots,1$
\begin{itemize}

\item ${\bmi R}_{ \uparrow \downarrow }^{\rho}=
{\bmi S}_{ \uparrow \downarrow }^{\rho}+{\bmi S}_{ \uparrow \uparrow 
}^{\rho}
{\bmi  E}_{\rho+1} {\bmi R}_{ \uparrow \downarrow }^{\rho+1} {\bmi  E}_{\rho+1} {\left({{\bmi  I}-{\bmi S}_{ \downarrow 
\uparrow }^{\rho}
{\bmi  E}_{\rho+1} {\bmi R}_{ \uparrow \downarrow }^{\rho+1} {\bmi  E}_{\rho+1} }\right)}^{-1}{\bmi S}_{ \downarrow 
\downarrow }^{\rho}
$

\item ${\bmi R}_{ \uparrow \uparrow }^{\rho}= {\bmi S}_{ \uparrow \uparrow 
}^{\rho}
{\left({{\bmi  I}-{\bmi  E}_{\rho+1} {\bmi R}_{ \uparrow \downarrow }^{\rho+1} {\bmi 
E}_{\rho+1} {\bmi S}_{ \downarrow 
\uparrow }^{\rho}}\right)}^{-1}$

\item ${\bmi R}_{ \downarrow \downarrow }^{\rho}={\left({{\bmi  I}-{\bmi S}_{ 
\downarrow \uparrow }^{\rho} {\bmi  E}_{\rho+1} {\bmi R}_{ \uparrow \downarrow 
}^{\rho+1} {\bmi  E}_{\rho+1} }
\right)}^{-1}{\bmi S}_{ \downarrow \downarrow }^{\rho} $

\item ${\bmi R}_{ \downarrow \uparrow }^{\rho} = 
{\left({{\bmi  I}-{\bmi S}_{ \downarrow \uparrow }^{\rho}
{\bmi  E}_{\rho+1} {\bmi R}_{ \uparrow \downarrow }^{\rho+1} {\bmi  E}_{\rho+1} }\right)}^{-1}
{\bmi S}_{ \downarrow \uparrow }^{\rho}$

\item $
\left[{\begin{array}{c}{\bmi \Sigma}^{\rho}_{\uparrow}\\ 
{\bmi \Sigma}^{\rho}_{\downarrow}
\end{array}}\right]=-i\,\left({\begin{array}{cc}
\sqrt{{\bmi E}_{\rho}}&{\bf 0}\\ 
{\bf 0}&\sqrt{{\bmi E}_{\rho}}
\end{array}}\right)\,\hat{\bmi  D}{}^{T}_{z,\rho}(-{\bf k}_{x},-{k}_{y})\left({\bmi I}_8\otimes
{\bf DFT}^{T}(-{\bf k}_{x})  \right)\hat{\bmi \Phi}\left[{\CB F}_{\rho}  \right]
$

\item  $\left[{\begin{array}{c}
{\bmi t}_{ \uparrow }^{\rho}\\ {\bmi t}_{ \downarrow }^{\rho+1}\end{array}}\right]=
\left[{\begin{array}{cc}
{\bmi R}_{ \uparrow \downarrow }^{\rho}&{\bmi R}_{ \uparrow \uparrow }^{\rho}\\ 
{\bmi R}_{ \downarrow \downarrow }^{\rho}&{\bmi R}_{ \downarrow \uparrow 
}^{\rho}\end{array}}\right]
\left[{\begin{array}{c}{\bmi \Sigma}^{\rho}_{\downarrow}\\ 
{\bmi \Sigma}^{\rho+1}_{\uparrow} + {\bmi  E}_{\rho+1}{  \bmi t}_{ \uparrow }^{\rho+1}
\end{array}}\right]$
\end{itemize}
\item {\bf  End}

\begin{itemize}
\item  (*{\em top of layer stack\/ }*) 
\end{itemize}
\item ${\bmi R}_{ \downarrow \uparrow }^{0} = 
{\left({{\bmi  I}-{\bmi S}_{ \downarrow \uparrow }^{[f,{0}]}
{\bmi  E}_{1} {\bmi R}_{ \uparrow \downarrow }^{1} {\bmi  E}_{1} }\right)}^{-1}
{\bmi S}_{ \downarrow \uparrow }^{[f,{0}]}$


\item  ${\bmi t}_{ \downarrow }^{1}={\bmi R}_{ \downarrow \uparrow }^{0}
({\bmi \Sigma}^{1}_{\uparrow}+{\bmi  E}_{\rho}{  \bmi t}_{ \uparrow }^{1})$

\item  ${\mit \Delta}{{\bmi V}}_{ \downarrow }^{1}={\bmi t}_{ \downarrow 
}^{1}$

\begin{itemize}
\item  (*{\em downsweep\/ }*) 
\end{itemize}
\item {\bf  For} $\rho=1,\cdots,L-1$
\begin{itemize}
\item  $
\left[{\begin{array}{c}
{\mit \Delta}{{\bmi V}}_{ \uparrow }^{\rho}\\ {\mit \Delta}{{\bmi V}}_{ 
\downarrow }^{\rho+1}\end{array}}\right]=\left[{\begin{array}{c}{  \bmi R}_{ 
\uparrow \downarrow }^{\rho}\\ {  \bmi R}_{ \downarrow \downarrow 
}^{\rho}\end{array}}\right] {\bmi  E}_{\rho} {  {\bmi V}}_{ \downarrow 
}^{\rho}+\left[{\begin{array}{c}{  \bmi t}_{ \uparrow }^{\rho}\\ {  
\bmi t}_{ \downarrow }^{\rho+1}\end{array}}\right] $

\item $
{\CB B}_{\rho}=\hat{\bmi \Phi}{}^{-1}\left[\left({\bmi I}_8\otimes 
{\bf DFT}({\bf k}_{x}) \right)\hat{\bmi  D}_{z,\rho}({\bf k}_{x},{k}_{y})
\left({\begin{array}{c}
\sqrt{{\bmi E}_{\rho}}{\mit \Delta}{{\bmi V}}_{ \downarrow }^{\rho}\\ 
\sqrt{{\bmi E}_{\rho}}{\mit \Delta}{{\bmi V}}_{ \uparrow }^{\rho}
\end{array}}\right)\right]
$

\end{itemize}


\item {\bf  End}





\begin{itemize}
\item  (*{\em bottom of layer stack\/ }*) 
\end{itemize}

\item  	$
{\mit \Delta}{{\bmi V}}_{ \uparrow }^{L}=
{\bmi R}_{ \uparrow \downarrow }^{L}
{\bmi  E}_{L} {{\bmi V}}_{ \downarrow }^{L} +{\bmi t}_{ \uparrow }^{L}$

\item $
{\CB B}_{L}=\hat{\bmi \Phi}{}^{-1}\left[\left({\bmi I}_8\otimes 
{\bf DFT}({\bf k}_{x}) \right)\hat{\bmi  D}_{z,L}({\bf k}_{x},{k}_{y})
\left({\begin{array}{c}
\sqrt{{\bmi E}_{L}}{\mit \Delta}{{\bmi V}}_{ \downarrow }^{L}\\ 
\sqrt{{\bmi E}_{L}}{\mit \Delta}{{\bmi V}}_{ \uparrow }^{L}
\end{array}}\right)\right]
$


\item{\bf End Algorithm}

\end{itemize}

\vspace{2mm}

\end{minipage}
}
\end{figure}

%%%%%%%%%%%%%%%%%%%%%%%%%%%%%%%%%%%%%%%%%%%%%%%%%%%%%%%%%%%%%%%%%%%%%%%%%%%%%%%%%%%%%%%%%%%
%%%%%%%%%%%%%%%%%%%%%%%%%%%%%%%%%%%%%%%%%%%%%%%%%%%%%%%%%%%%%%%%%%%%%%%%%%%%%%%%%%%%%%%%%%%
%%%%%%%%%%%%%%%%%%%%%%%%%%%%%%%%%%%%%%%%%%%%%%%%%%%%%%%%%%%%%%%%%%%%%%%%%%%%%%%%%%%%%%%%%%%
%%%%%%%%%%%%%%%%%%%%%%%%%%%%%%%%%%%%%%%%%%%%%%%%%%%%%%%%%%%%%%%%%%%%%%%%%%%%%%%%%%%%%%%%%%%


